\providecommand{\spnHS}{\mathrm{HS}}
\providecommand{\spncodim}{\operatorname{codim}}
\providecommand{\spnim}{\operatorname{im}}
\providecommand{\spnE}{\mathbb{E}}
\providecommand{\spnP}{\mathbb{P}}
\providecommand{\spnTV}{d_{\mathrm{TV}}}
\providecommand{\spnStab}{\operatorname{Stab}}
\providecommand{\spnMat}{\operatorname{Mat}}
\providecommand{\spnDD}{\mathfrak{D}}
\providecommand{\spnT}{\mathsf{T}}
\providecommand{\spnqbin}[3]{\genfrac{[}{]}{0pt}{}{#1}{#2}_{#3}}
\providecommand{\spnOrth}{\mathcal{O}}
\providecommand{\spnTame}{\mathcal{T}}
\providecommand{\spnEmb}{\mathcal{E}}
\providecommand{\spnAlt}{\operatorname{Alt}}
\providecommand{\spncoker}{\operatorname{coker}}
\providecommand{\torsHs}{H_{*}}
\providecommand{\torslen}{\operatorname{length}}
\providecommand{\torscoker}{\operatorname{coker}}
\providecommand{\torsadj}{\operatorname{adj}}
\providecommand{\torsFitt}{\operatorname{Fitt}}
\providecommand{\torsZ}[2]{Z_{[#1,#2]}}
\providecommand{\torsU}{\mathcal{U}}
\providecommand{\torsfiles}[1]{\begin{itemize}[nosep,leftmargin=1.5em,label=--]\footnotesize #1\end{itemize}}
\renewcommand{\arraystretch}{1.25}
\section{Torsion specialization: exact costs, obstructions and coefficient fields}\label{sec:cost}
% Local helper: a compact list of delivered files (paths are too long for running text).
\providecommand{\costfiles}[1]{\begin{itemize}[nosep,leftmargin=1.5em,label=--]\footnotesize #1\end{itemize}}

The selection theorems of \cref{sec:tors}, in particular \cref{tors:thm:selection}, produce a
root of unity $\zeta$ of $p$-power order such that an integral $p$-adic polynomial or power
series of bounded Gauss valuation does not gain valuation at $\zeta-1$. This section treats
the opposite side of the same problem: how much valuation budget is needed to defeat every
such specialization on a prescribed set of primitive shells. The answer depends on three
choices: whether one polynomial is compared with one fixed Gauss valuation or several
polynomials are compared individually; whether a positive tolerance is allowed; and which
coefficient field is admitted. Each case is settled by an exact formula or by an exact finite
computation; for packets with a prescribed number of members, which interpolate between the
first two cases, the cost is also determined asymptotically (\cref{cost:sec:cardinality}).
All results are local statements about $p$-adic polynomials. None of them uses a
Riemann hypothesis.

\subsection{Separate exact bounds over a fixed coefficient field}\label{tors:sec:fixedfield}
\sources{grh_fixed_field_transfer_2026-11-25; grh-torsion-and-prime-sampling}

Let $K/\Qp$ be a finite extension inside $\Cp$, with ring of integers $\OO_K$, absolute
ramification index $e=e(K/\Qp)$, and
\[
  w=\vp(e) .
\]
When the coefficients are restricted to $\OO_K$, the value group forces multiplicities of
roots on a shell. This makes separate comparisons with tolerance zero possible for
finite families of any cardinality.

\begin{lemma}[Multiplicity at one torsion radius]\label{tors:lem:residual}
Fix $h\ge1$ and $z_0\in Z_h$. For a nonzero $g=\sum_jc_jT^j\in\OO_K[T]$ choose $b$ in a
finite extension of $K$ with $v(b)=G_{r_h}(g)$, and let
\[
  \mathcal I_g(Y)=\overline{b^{-1}g(z_0Y)}
\]
be the reduction in the residue field of a common finite extension.
\begin{enumerate}[label=\textup{(\alph*)}]
\item $\mathcal I_g\ne0$, and the number of its nonzero roots, counted with multiplicity,
  is $m_h(g)$.
\item Every nonzero root of $\mathcal I_g$ has multiplicity at least
  $P_h:=p^{\max\{h-1-w,\,0\}}$.
\item If $z\in Z_h$ and $v(g(z))>t$ for some real $t\ge G_\rho(g)$ with $\rho\ge r_h$,
  then $\overline{z/z_0}$ is a nonzero root of $\mathcal I_g$.
\end{enumerate}
\end{lemma}

\begin{proof}
(a) Factor $g$ over $\Cp$. After the scaling $T=z_0Y$ and division by a constant of the
right valuation, a root of valuation smaller than $r_h$ contributes a nonzero constant
to the reduction, a root of valuation larger than $r_h$ contributes the factor $Y$, and
a root $\alpha$ of valuation exactly $r_h$ contributes the factor
$Y-\overline{\alpha/z_0}$ with a nonzero root.

(b) A coefficient $c_j$ survives the reduction exactly when
$v(c_j)+j\,r_h=G_{r_h}(g)$. If the indices $j$ and $k$ both survive, then
$(j-k)r_h\in e^{-1}\Z$, hence $e_h/\gcd(e,e_h)$ divides $j-k$. The $p$-part of this
divisor is $P_h$. After removing a power of $Y$, the polynomial $\mathcal I_g$ is
therefore a polynomial in $Y^{P_h}$, and over an algebraic closure of the residue field,
which has characteristic $p$, it is a $P_h$-th power.

(c) Monotonicity gives $t\ge G_\rho(g)\ge G_{r_h}(g)$, so $b^{-1}g(z)$ has positive
valuation. The quotient $z/z_0$ is a unit, and its residue is a root of $\mathcal I_g$.
\end{proof}

\begin{theorem}[Separate exact bounds over a fixed field; \stU]\label{tors:thm:fixedfield}
Let $S$ be a nonempty finite set of positive integers and let $\rho\ge\max_{h\in S}r_h$.
Let $g_1,\dots,g_n\in\OO_K[T]$ be nonzero, with real targets $t_i\ge G_\rho(g_i)$. Put
\[
  a_h=p^{-\min\{w,\,h-1\}},\qquad C_K(S)=\sum_{h\in S}a_h .
\]
\begin{enumerate}[label=\textup{(\alph*)}]
\item If $h\in S$ and every $z\in Z_h$ has $v(g_i(z))>t_i$ for at least one $i$, then
  $r_h\sum_im_h(g_i)\ge a_h$.
\item If $\sum_it_i<C_K(S)$, there are $h\in S$ and $z\in Z_h$ with
  \[
    v(g_i(z))\le t_i\qquad\text{for every }i .
  \]
\item Let $h_0\ge1$ satisfy $r_{h_0}\le\rho$ and let $\sum_it_i\le B$. Then the
  conclusion of \textup{(b)} holds with some $h$ in $[h_0,H]$, where $H$ is the first
  integer $\ge h_0$ with $\sum_{h=h_0}^{H}p^{-\min\{w,h-1\}}>B$. Always
  $H\le h_0+\floor{p^wB}$, with equality when $h_0\ge w+1$.
\end{enumerate}
Neither $n$ nor the degrees are bounded. Witnesses $x_i$ of depth $\rho$ with
$t_i=v(g_i(x_i))$ give a special case; the witnesses need not lie in $K$.
\end{theorem}

\begin{proof}
(a) Write $z_0=\zeta_0-1$ with $\zeta_0$ primitive of order $p^h$. For $a=1,\dots,p-1$
the points $z_a=\zeta_0^a-1$ belong to $Z_h$ and
\[
  \overline{z_a/z_0}=\overline{1+\zeta_0+\dots+\zeta_0^{a-1}}=a\in\F_p^\times .
\]
By \cref{tors:lem:residual}(c) each of these $p-1$ residues is a root of
$\mathcal I_{g_i}$ for at least one $i$, with multiplicity at least $P_h$ by
\cref{tors:lem:residual}(b). With \cref{tors:lem:residual}(a) this gives
$\sum_im_h(g_i)\ge(p-1)P_h$, and
\[
  r_h\sum_im_h(g_i)\ge\frac{(p-1)\,p^{\max\{h-1-w,0\}}}{(p-1)\,p^{h-1}}
  =p^{-\min\{w,h-1\}}=a_h .
\]

(b) If the conclusion failed for every $h\in S$, then by (a) and
\cref{tors:lem:shells}(a)
\[
  C_K(S)=\sum_{h\in S}a_h\le\sum_i\sum_{h\in S}r_h\,m_h(g_i)\le\sum_iG_\rho(g_i)
  \le\sum_it_i .
\]
The roots of the $g_i$ are not assumed to be roots of unity.

(c) Apply (b) to $S=[h_0,H]\cap\Z$. Every $a_h$ is at least $p^{-w}$, so
$\floor{p^wB}+1$ consecutive levels have total cost greater than $B$; if $h_0\ge w+1$,
every $a_h$ with $h\ge h_0$ equals $p^{-w}$.
\end{proof}

\begin{remark}\label{tors:rem:fixedfield}
(1) The number $C_K(S)$ is exactly the least budget of a family over $\OO_K$ for which the
conclusion of \cref{tors:thm:fixedfield}(b) fails; it is attained by one explicit
binomial for each $h\in S$, and equality $\sum_it_i=C_K(S)$ already permits failure. In
particular the endpoint $H$ in \textup{(c)} is optimal. These facts are proved in
\cref{cost:thm:field,cost:cor:fieldendpoint}.

(2) The cost depends on $K$ only through the wild part $p^w$ of $e$. Every tamely
ramified $K$ has the endpoint $h_0+\floor{B}$ of $\Qp$. The earlier fixed-field count,
which uses only that a shell carrying a root costs at least $1/e$, gives
$h_0+\floor{eB}$; for $p=2$, $h_0=1$, $B=3$ this is $316$ for $e=105$ and $25$ for $e=8$,
against $4$ and $14$ from \cref{tors:thm:fixedfield}(c); further values are in
\artifact{grh_fixed_field_transfer_2026-11-25/data/bound_comparison.csv}.

(3) For one polynomial over $\OO_K$ the separate and the aggregate comparison coincide,
and both \cref{tors:thm:selection} and \cref{tors:thm:fixedfield}(c) apply. When
$w\le h_0-1$ the fixed-field endpoint is $h_0+\floor{p^wB}$, which is at most the
universal endpoint $h_0+\floor{p^{h_0-1}B}$.
\end{remark}

\begin{corollary}[A weighted sampling guarantee; \stU]\label{tors:cor:sampling}
In the situation of \cref{tors:thm:fixedfield}, let $\sum_it_i\le B$. Choose $h\in S$
with probability $a_h/C_K(S)$ and then $z$ uniformly in $Z_h$. The probability that
$v(g_i(z))\le t_i$ for every $i$ is at least $\max\{0,\,1-B/C_K(S)\}$.
\end{corollary}

\begin{proof}
Put $b_h=r_h\sum_im_h(g_i)$. By \cref{tors:lem:residual} at most $\sum_im_h(g_i)/P_h$ of
the $p-1$ residues in $\F_p^\times$ are roots of some $\mathcal I_{g_i}$, and every
residue class contains $p^{h-1}$ points of $Z_h$. The proportion of failing points in
$Z_h$ is therefore at most $\min\{1,b_h/a_h\}$, and the weighted probability of failure is
at most $\sum_{h\in S}b_h/C_K(S)\le B/C_K(S)$ by \cref{tors:lem:shells}(a).
\end{proof}

\paragraph{Several variables.}
Choosing the coordinates one after another, as in the proof of
\cref{tors:thm:selection}(b), enlarges the coefficient field by a cyclotomic extension at
every step, and the resulting exponents grow like an iterated exponential
(\cref{tors:rem:tower}). The following reduction avoids this. It restricts the family to
one of finitely many multiplicative curves, chosen before any torsion parameter and
defined by integer coefficients, so that the coefficient field stays fixed until a single
final specialization in one variable.

\begin{lemma}[Finitely many slopes preserve separate bounds]\label{tors:lem:slopes}
Let $(F_i)_{i\in I}$ be a finite family of nonzero integral polynomials in the variables
$X,Y,V_1,\dots,V_k$ over $\OO$. Fix $\rho>0$, put $A_i=G_\rho(F_i)$, and assume
$\sum_iA_i\le B$. Set
\[
  M=\floor*{\frac B\rho}+1,\qquad r=\frac{\rho}{M^2},\qquad b_n(T)=(1+T)^n-1 .
\]
Then for some integer $n$ with $1\le n\le M$, simultaneously for all $i$,
\[
  G_{(r,\rho,\dots,\rho)}\bigl(F_i(T,b_n(T),V_1,\dots,V_k)\bigr)\le A_i ,
\]
where the weight $r$ belongs to $T$. The same holds after lowering the weights of the
$V_j$ to $r$. No bound on the degrees or on $\#I$ is required.
\end{lemma}

\begin{proof}
Consider one $F$ with $A=G_\rho(F)$ and put $l=\floor{A/\rho}+1$. Call a slope $n$
\emph{bad} if the Gauss valuation of $F(T,b_n(T),V)$ with weights $(r,\rho,\dots,\rho)$
is strictly greater than $A$; this includes the case that the restriction is zero. We
show that fewer than $l$ slopes in $\set{1,\dots,M}$ are bad. Suppose that a set
$J\subseteq\set{1,\dots,M}$ of $l$ slopes is bad, and divide $F$ in $Y$ by the monic
polynomial
\[
  Q_J(X,Y)=\prod_{n\in J}\bigl(Y-b_n(X)\bigr).
\]
Quotient and remainder are integral: $F=Q_JS+R$ with $\deg_YR<l$. Each factor of $Q_J$
has Gauss valuation exactly $\rho$ at radius $\rho$, so $G_\rho(Q_JS)\ge l\rho>A$ and
$G_\rho(R)=A$.

Choose a monomial of $R$ attaining this minimum. If its degree in $X,Y$ is zero, then
$G_\rho(R(0,0,V))\le A$. But $R(0,0,V)=F(0,0,V)$ is the constant coefficient in $T$ of
every restriction $F(T,b_n(T),V)$, so every slope would be good, a contradiction.
Otherwise the degree in $X,Y$ is at least one, and lowering the weights of $X$ and $Y$
from $\rho$ to $r$ gives
\begin{equation}\label{tors:eq:slopesaving}
  G_{(r,r,\rho,\dots,\rho)}(R)\le A-(\rho-r).
\end{equation}
Lagrange interpolation in $Y$ at the $l$ distinct nodes $b_n(X)$ gives
\begin{gather*}
  R(X,Y,V)=\sum_{n\in J}F(X,b_n(X),V)\,
  \frac{\prod_{m\in J\setminus\{n\}}(Y-b_m(X))}{D_n(X)},\\
  D_n(X)=\prod_{m\in J\setminus\{n\}}\bigl(b_n(X)-b_m(X)\bigr).
\end{gather*}
The Gauss valuation with weights $(r,r,\rho,\dots,\rho)$ is multiplicative and extends to
the field of rational functions. Each numerator product has Gauss valuation $(l-1)r$.
For a positive integer $d$ the coefficient of $T^{p^{\vp(d)}}$ in $(1+T)^d-1$ is a
$p$-adic unit, because
$(1+T)^d\equiv(1+T^{p^{\vp(d)}})^{d/p^{\vp(d)}}\pmod p$. Consequently
\[
  G_r(b_n-b_m)\le p^{\vp(n-m)}\,r\qquad(n\ne m).
\]
Define
\[
  C_J=\max_{n\in J}\sum_{m\in J\setminus\{n\}}\bigl(p^{\vp(n-m)}-1\bigr).
\]
Since all slopes in $J$ are bad, the interpolation identity yields
\begin{equation}\label{tors:eq:slopeinterp}
  G_{(r,r,\rho,\dots,\rho)}(R)>A-r\,C_J .
\end{equation}
For $M\ge2$ one has $p^{\vp(n-m)}\le\abs{n-m}\le M-1$ and $l\le M$, so
$C_J\le(l-1)(M-2)\le(M-1)(M-2)$; for $M=1$, $C_J=0$. Thus $C_J+1\le M^2$ and
$r\,C_J\le\rho-r$, so \eqref{tors:eq:slopesaving} and \eqref{tors:eq:slopeinterp}
contradict each other.

Each $F_i$ therefore has at most $\floor{A_i/\rho}$ bad slopes. The union of the bad sets
has at most $\sum_i\floor{A_i/\rho}\le\floor{B/\rho}=M-1$ elements, so at least one of the
$M$ slopes is good for all $i$.
\end{proof}

\begin{theorem}[Several variables without a ramification tower; \stU]\label{tors:thm:notower}
Let $(F_i)_{i\in I}$ be a finite family of nonzero polynomials in
$\OO_K[T_1,\dots,T_d]$, $d\ge1$. Fix $\rho>0$, put $A_i=G_\rho(F_i)$, and assume
$\sum_iA_i\le B$. Let $h_0\ge1$ satisfy $r_{h_0}\le\rho$. There are roots of unity
$\zeta_1,\dots,\zeta_d$ of orders $p^{h_1},\dots,p^{h_d}$ with
\[
  h_0\le h_j\le H,\qquad v\bigl(F_i(\zeta_1-1,\dots,\zeta_d-1)\bigr)\le A_i
  \quad\text{for every }i,
\]
where
\begin{equation}\label{tors:eq:notower}
  H=h_0+\ceil*{(3^{d-1}-1)\log_p\Bigl(1+\frac B\rho\Bigr)}+\floor{p^wB} .
\end{equation}
\end{theorem}

\begin{proof}
Put $\rho_0=\rho$ and, for $j=0,\dots,d-2$,
\[
  M_j=\floor*{\frac{B}{\rho_j}}+1,\qquad \rho_{j+1}=\frac{\rho_j}{M_j^2}.
\]
At stage $j$ apply \cref{tors:lem:slopes}, with radius $\rho_j$, to the parameter $T$ of
the curve constructed so far (initially $T=T_1$) and the next variable $T_{j+2}$, and
then lower the weights of the remaining variables to $\rho_{j+1}$. Every separate Gauss
valuation decreases or stays the same, and all coefficients remain in $\OO_K$. We obtain
integers $a_1=1,a_2,\dots,a_d\ge1$ with $a_{j+2}\le M_j$ and
\[
  P_i(T)=F_i\bigl((1+T)^{a_1}-1,\dots,(1+T)^{a_d}-1\bigr)\in\OO_K[T],\qquad
  G_{\rho_{d-1}}(P_i)\le A_i .
\]
Write $\Pi=\prod_{j=0}^{d-2}M_j$, an empty product being $1$. Then $a_j\le\Pi$ and
$\rho_{d-1}=\rho/\Pi^2$.

Put $\bar h=h_0+\ceil{2\log_p\Pi}$, so that $r_{\bar h}\le\rho/\Pi^2$. Apply
\cref{tors:thm:fixedfield}(b) to the $P_i$ with the radius $\rho_{d-1}$, the targets
$A_i\ge G_{\rho_{d-1}}(P_i)$, and $S=[\bar h,\bar h+\floor{p^wB}]\cap\Z$: every level
costs at least $p^{-w}$, so $C_K(S)>B$. This gives a root of unity $\zeta$ of order
$p^h$ with
\[
  \bar h\le h\le\bar h+\floor{p^wB},\qquad v(P_i(\zeta-1))\le A_i\quad\text{for all }i .
\]
The targets dominate the Gauss valuations of the $P_i$; they need not be values at a
common witness. Set $\zeta_j=\zeta^{a_j}$. Its order exponent is $h-\vp(a_j)$, which is
at most $h$ and at least
$h_0+\ceil{2\log_p\Pi}-\floor{\log_p\Pi}\ge h_0$.

For the explicit bound put $u_j=1+B/\rho_j$. Then $M_j\le u_j$ and
\[
  u_{j+1}=1+(u_j-1)M_j^2\le u_j^3,\qquad u_j\le u_0^{3^j},\qquad
  2\log_p\Pi\le(3^{d-1}-1)\log_p\Bigl(1+\frac B\rho\Bigr).
\]
This gives \eqref{tors:eq:notower}. The cases $d=1$, $B=0$ and $M_j=1$ are included.
\end{proof}

\begin{corollary}[Bounded series and arbitrary families over $\OO_K$; \stU]
\label{tors:cor:notowerseries}
\Cref{tors:thm:notower} holds, with the same $H$, for nonzero series
$F_i\in\OO_K[[T_1,\dots,T_d]]$, and for an arbitrary family of such series for which
$\sum_{i\in J}G_\rho(F_i)\le B$ for every finite subfamily $J$.
\end{corollary}

\begin{proof}
For a finite family apply \cref{tors:lem:truncation} with $a=0$, $M=B$ and $r=r_H$. Since
$\rho\ge r_{h_0}\ge r_H$, the truncations have the Gauss valuations $A_i\le B$ of the
$F_i$, and the omitted tails have valuation greater than $B$ at every point of the grid.
The polynomial theorem gives a tuple with $v(P_i(z))\le A_i\le B$, hence
$v(F_i(z))=v(P_i(z))$. Since $H$ depends only on $p,w,d,\rho,B,h_0$, the truncation
degree is chosen without circularity. For an arbitrary family the sets of good tuples of
the finite subfamilies are nonempty subsets of the finite grid $\torsZ{h_0}{H}^d$ with
the finite intersection property.
\end{proof}

\begin{remark}[Comparison with iterated specialization]\label{tors:rem:tower}
Specializing one coordinate at a time with a fixed-field count in one variable gives
exponents $h_j\le M_j$ with the recurrence
\[
  M_1=h_0+\floor{eB},\qquad M_{j+1}=h_0+\floor{e\,\varphi(p^{M_j})\,B},
\]
because after $j$ steps the coefficients lie in $K(\mu_{p^{M_j}})$. For $p=2$, $K=\Q_2$,
$\rho=1$, $B=3$, $h_0=1$ the bound \eqref{tors:eq:notower} is $H=2\cdot3^{d-1}+2$:
\begin{center}
\begin{tabular}{@{}rrr@{}}
\toprule
$d$ & recurrence & \eqref{tors:eq:notower}\\
\midrule
1 & $4$ & $4$\\
2 & $25$ & $8$\\
3 & $50{,}331{,}649$ & $20$\\
4 & $1+3\cdot2^{50{,}331{,}648}$ & $56$\\
\bottomrule
\end{tabular}
\end{center}
These are sufficient bounds. The cost in one variable is sharp; the factor $3^{d-1}$ and
the range of slopes come from a convenient interpolation estimate and are not claimed to
be optimal. The question whether the iterated growth is necessary is thereby answered in
the negative.
\end{remark}

\begin{example}[One slope is not enough]\label{tors:ex:slopes}
A single prescribed multiplicative curve cannot serve in \cref{tors:thm:notower}: for
coprime positive integers $a,b$ the curve $((1+T)^a-1,(1+T)^b-1)$ is annihilated by
$(1+T_1)^b-(1+T_2)^a$, an integral polynomial of Gauss valuation exactly $\rho$, since at
least one of its linear coefficients is a $p$-adic unit. Concretely, over $\Z_2[X,Y]$ let
\[
  F_n(X,Y)=Y-\bigl((1+X)^n-1\bigr),\qquad n=1,2,3 .
\]
At $\rho=1$ all three Gauss valuations are $1$, and $(X,Y)=(0,2)$ is a common witness at
which all three valuations equal $1$; the budget is $3$. Each substitution
$Y=(1+X)^n-1$, $n=1,2,3$, makes one member of the family identically zero.
\Cref{tors:lem:slopes} considers $M=4$ slopes and selects the fourth. If $\zeta$ is a
primitive eighth root of unity, the tuple $X=\zeta-1$, $Y=\zeta^4-1=-2$ gives
\[
  v(\zeta^4-\zeta)=\tfrac14,\qquad v(\zeta^4-\zeta^2)=\tfrac12,\qquad
  v(\zeta^4-\zeta^3)=\tfrac14,
\]
all at most $1$. The theorem allows $H=8$; this example needs only $H=3$. The exact
check of this family is recorded in:
\torsfiles{
\item \artifact{grh_fixed_field_transfer_2026-11-25/checks/finite_slope_checks.json}}
\end{example}

\begin{remark}[Order of quantifiers]\label{tors:rem:quantifiers}
In \cref{tors:thm:fixedfield,tors:thm:notower} the data $K,p,d,\rho,B$ are fixed first,
then $H$ is chosen, and then every admissible family over $\OO_K$ is handled, with
arbitrary degrees and arbitrarily many members. The coefficients must lie in the fixed
field, or in fields with a known uniform bound for the wild ramification exponent $w$.
By contrast \cref{tors:thm:uniform} handles arbitrary integral coefficients after $H$ is
chosen, at positive tolerance. A finite family always lies in some finite extension, but
choosing that extension after inspecting the family does not give a uniform bound in the
former order of quantifiers. The fixed-field theorems can therefore replace the
tolerance in interpolation settings over a fixed field; they do not remove the tolerance
from every use of the selector in the source.
\end{remark}

\subsection{Packets over a fixed coefficient field: reduction and zero tolerance}\label{cost:sec:fields}
{\raggedright\sources{grh_fixed_field_transfer_2026-11-25; grh_sharp_torsion_and_symplectic_blocks; grh-torsion-and-prime-sampling; grh-torsion-and-zero-data; grh_method_transfers_sharp_torsion}}

The extremal packets of Theorem~\ref{cost:thm:packet} live in cyclotomic fields whose ramification
grows with the largest forbidden shell. From now on the coefficient field is prescribed:
$K\subset\Cp$ is a finite extension of $\Qp$ with ring of integers $\OO_K$, absolute
ramification index $e_K$, wild exponent $w=v_p(e_K)$ and residue field $k$. We stay in
setting (C). The roots of valuation $r_h$ of a polynomial over $K$ form a Galois-stable
multiset, and this is the only new constraint.

\begin{proposition}[Reduction to a monic test on one shell; \stU]\label{cost:prop:monic}
For $h\ge1$ and $\delta\ge0$ let $d_{K,h}(\delta)$ be the least degree $m$ of a monic
polynomial $P\in\OO_K[T]$ with
\begin{equation}\label{cost:eq:monictest}
 v(P(z))>m\,r_h+\delta\qquad\text{for every }z\in Z_h,
\end{equation}
and put $c_{K,h}(\delta)=d_{K,h}(\delta)\,r_h$.
\begin{enumerate}[label=\textup{(\roman*)}]
\item $1\le d_{K,h}(\delta)\le e_h$. A monic polynomial of least degree satisfying
\eqref{cost:eq:monictest} has all its roots of valuation $r_h$, and
$c_{K,h}(\delta)=v(P(0))\in v(K^\times)=e_K^{-1}\Z$.
\item For every nonempty finite $S$ with $\rho\ge\max_{h\in S}r_h$,
\[
 C_K(S,\delta)=\sum_{h\in S}c_{K,h}(\delta),
\]
and the minimum is attained by one polynomial per shell, each with witness $0$.
\end{enumerate}
\end{proposition}

\begin{proof}
(i) The polynomial $\Phi_{p^h}(1+T)$ vanishes on $Z_h$ and has degree $e_h$. Let $P$ satisfy
\eqref{cost:eq:monictest} and split $P=P_{\mathrm{on}}P_{\mathrm{off}}$ according to
whether a root has valuation $r_h$ or not. Both factors are monic, integral and defined
over $K$, because the valuation is Galois invariant. A root $\beta$ of $P_{\mathrm{off}}$
has $v(z-\beta)=\min\set{r_h,v(\beta)}\le r_h$, so $P_{\mathrm{on}}$ satisfies
\eqref{cost:eq:monictest} with its own degree. Hence a polynomial of least degree has all
roots on the radius, and then $m\,r_h=v(P(0))$.

(ii) Let a packet over $\OO_K$ block $S$. For $h\in S$ let $P_h$ be the monic polynomial
whose roots are the roots of valuation $r_h$ of all the $g_i$; it lies in $\OO_K[T]$. By
\eqref{cost:eq:charge}, $\mathcal L(z)\le v(P_h(z))-(\deg P_h)r_h$ on $Z_h$, so $P_h$
satisfies \eqref{cost:eq:monictest} and $\deg P_h\ge d_{K,h}(\delta)$; the second part of
\eqref{cost:eq:charge} gives the lower bound. Conversely let $P_h$ be a polynomial of least
degree in (i), used as a packet member with witness $0$ and threshold
$v(P_h(0))=G_\rho(P_h)$. On its own shell its loss exceeds $\delta$. On another shell
$Z_{h'}$ its valuation is $\deg P_h\cdot\min\set{r_h,r_{h'}}$, which is at most its
threshold, so it contributes no loss. The costs add.
\end{proof}

Since $\OO_K\subset\OO$, one has $c_{K,h}(\delta)\ge c_h(\delta)$, with equality when $K$
contains $\zeta_{p^h}$, by Theorem~\ref{cost:thm:packet}. The function $c_{K,h}$ is nondecreasing
in $\delta$ and takes values in $e_K^{-1}\Z\cap(0,1]$.

\begin{theorem}[Exact price at zero tolerance over a finite coefficient field; \stU]\label{cost:thm:field}
For every finite extension $K/\Qp$ and every $h\ge1$,
\[
 c_{K,h}(0)=p^{-\min\set{w,\,h-1}},\qquad\text{hence}\qquad
 C_K(S,0)=\sum_{h\in S}p^{-\min\set{w,\,h-1}}
\]
for every nonempty finite $S$ with $\rho\ge\max_{h\in S}r_h$; this is the number $C_K(S)$
of \cref{tors:thm:fixedfield}. The minimum is attained by one binomial per shell with
common witness $0$. In particular the strict inequality $\sum_it_i<C_K(S)$ in
\cref{tors:thm:fixedfield}\textup{(b)} cannot be replaced by equality.
\end{theorem}

\begin{proof}
Fix $h$ and abbreviate $n=e_h$, $r=r_h$ and $s=\min\set{w,h-1}$.

\emph{Lower bound.} Let $P\in\OO_K[T]$ be monic of degree $m=d_{K,h}(0)$ with
$v(P(z))>mr$ on $Z_h$. By Proposition~\ref{cost:prop:monic}(i) all its roots have
valuation $r$, so $G_r(P)=mr$. The packet consisting of $P$ with threshold $mr$ fails at
every point of $Z_h$, and \cref{tors:thm:fixedfield}\textup{(a)}, with $S=\set h$ and
radius $r$, gives $mr\ge p^{-s}$. Thus $c_{K,h}(0)\ge p^{-s}$.

\emph{The binomial.} Put $m=n/p^s$. The polynomial $\Phi_{p^h}(1+T)$ is Eisenstein with
constant coefficient $p$, so $\overline{z^n/p}=-1$ for every $z\in Z_h$. There is $c\in K$
with $v(c)=p^{-s}$ and $\overline{c^{p^s}/p}=-1$: start with a uniformizer to the power
$e_K/p^s$ and adjust by a unit, using that the $p^s$-th power map is bijective on $k$. For
every $z\in Z_h$ the unit $z^m/c$ reduces to $1$, since its $p^s$-th power does. All roots
of $T^m-c$ have valuation $r$, and
\[
 v(z^m-c)>p^{-s}=v(c)=m\,r.
\]
So $T^m-c$ passes the test at tolerance zero with $m\,r=p^{-s}$. The formula for
$C_K(S,0)$ follows from Proposition~\ref{cost:prop:monic}(ii), and at cost exactly
$C_K(S,0)$ the binomials block $S$.
\end{proof}

The price depends only on the wild part $p^w$ of the ramification index. The lower bound
does not assume that the roots of the packet are themselves roots of unity.

\begin{corollary}[Tame and unramified fields, and exact endpoints; \stU]\label{cost:cor:fieldendpoint}
\leavevmode
\begin{enumerate}[label=\textup{(\roman*)}]
\item If $K/\Qp$ is tamely ramified, then $c_{K,h}(0)=1$ for all $h$, and
$C_K(S,0)=\abs S$: at tolerance zero every shell costs one full unit, as over $\Qp$, and
the cost is attained by the polynomials $\Phi_{p^h}(1+T)$. If $K/\Qp$ is unramified, the
same holds at every tolerance $\delta\ge0$: $c_{K,h}(\delta)=1$ and
$C_K(S,\delta)=\abs S$.
\item For $S=\set{h_0,\ldots,H}$, the endpoint of
\cref{tors:thm:fixedfield}\textup{(c)}, namely the first $H\ge h_0$ with
$\sum_{h=h_0}^Hp^{-\min\set{w,h-1}}>B$, is the least endpoint
guaranteeing, for every packet over $\OO_K$ of cost at most $B$, a point with
$v(g_i(z))\le t_i$ for all $i$. For $h_0\ge w+1$ it equals $h_0+\floor{p^wB}$.
\item Let $K/\Qp$ be unramified. The least value of $G_\rho(g)$ among nonzero
$g\in\OO_K[T]$ blocking $Z_{h_0},\ldots,Z_H$ at tolerance zero in the aggregate sense of
Definition~\ref{cost:def:aggregate} is $H-h_0+1$, attained by
$\prod_{h=h_0}^H\Phi_{p^h}(1+T)\in\Z[T]$. Thus for such coefficients the least universal
endpoint for $G_\rho(g)\le B$ is $h_0+\floor B$, in place of $h_0+\floor{p^{h_0-1}B}$ for
unrestricted coefficients.
\end{enumerate}
\end{corollary}

\begin{proof}
(i) The tame case is Theorem~\ref{cost:thm:field} with $w=0$. In the unramified case
$\Phi_{p^h}(1+T)$ vanishes throughout shell $h$ and has value $p$ at $T=0$, so one such
polynomial per shell blocks $S$ at every tolerance at total cost $\abs S$, and the
zero-tolerance value $\abs S$ is a lower bound at every tolerance. (ii) Sufficiency is
\cref{tors:thm:fixedfield}\textup{(c)}. If $h_0\le H'$ and
$\sum_{h=h_0}^{H'}p^{-\min\set{w,h-1}}\le B$, the binomials of Theorem~\ref{cost:thm:field}
for the shells $h_0,\ldots,H'$ form a packet of cost at most $B$ that blocks all of them,
so no endpoint below the stated one suffices. (iii) A single polynomial with threshold
$G_\rho(g)$ is a packet with one member, so a blocker has $G_\rho(g)\ge C_K(S,0)=\abs S$,
by Theorem~\ref{cost:thm:field} with $w=0$; the product of the shifted cyclotomic
polynomials vanishes on the whole window and has $G_\rho=H-h_0+1$ by
Lemma~\ref{cost:lem:roots}.
\end{proof}

A numerical comparison of the endpoint in (ii) with the earlier bound $h_0+\floor{e_KB}$ of
Proposition~\ref{cost:prop:wholeshell} is given in \cref{tors:rem:fixedfield}. Exact rational
checks of the zero-tolerance formula are in:
\costfiles{
\item \artifact{grh_fixed_field_transfer_2026-11-25/checks/zero_tolerance_checks.json}}

Blocking requires a bad point on every shell. A shell with no root at its radius has no bad
point at all, and such a shell exists as soon as the budget is exhausted by the
ramification.

\begin{proposition}[Whole shells of zero loss; \stU]\label{cost:prop:wholeshell}
Let $(g_i,t_i)_i$ be a packet over $\OO_K$ and $S$ a finite set of positive integers with
$\rho\ge\max_{h\in S}r_h$. Let $S'\subseteq S$ be the set of $h$ such that some $g_i$ has a
root of valuation $r_h$. Then
\[
 \sum_{h\in S'}\frac1{\gcd(e_K,e_h)}\le\sum_i\Delta_\rho(g_i)\le\sum_it_i,
\]
and for $h\in S\setminus S'$ every $z\in Z_h$ satisfies
$v(g_i(z))=G_{r_h}(g_i)\le t_i$ for all $i$ simultaneously. In particular, among any
$\floor{e_KB}+1$ shells of radius at most $\rho$ there is one on which a packet of cost at
most $B$ has no loss at any point. For unramified $K$, $\floor B+1$ shells suffice, and this
number is optimal at every tolerance.
\end{proposition}

\begin{proof}
An irreducible factor over $K$ of some $g_i$ that has a root of valuation $r_h$ has all its
roots of that valuation, by invariance of the extended valuation. If its degree is $D$, its
norm has valuation $D/e_h\in e_K^{-1}\Z$, so $D$ is a multiple of $e_h/\gcd(e_K,e_h)$. Hence a
shell in $S'$ carries at least $e_h/\gcd(e_K,e_h)$ roots and costs at least
$1/\gcd(e_K,e_h)\ge1/e_K$ in \eqref{cost:eq:budget}. On a shell without roots at its radius,
\eqref{cost:eq:shellvalue} gives $v(g_i(z))=G_{r_h}(g_i)$. Optimality for unramified $K$ is
Corollary~\ref{cost:cor:fieldendpoint}(i).
\end{proof}

For positive tolerance the function $c_{K,h}$ is an invariant of $K$ that is not determined
by the ramification index (Remark~\ref{cost:rem:twofields}), but it is always computable.

\begin{proposition}[A finite computation at every tolerance; \stU]\label{cost:prop:mod}
Fix $h$ and $\delta\ge0$. Let $L=K(\xi)$ for a primitive $p^h$-th root of unity $\xi$,
$E=e(L/\Qp)$, and $\varpi_L$ a uniformizer of $L$. For $m=1,\ldots,e_h$ put
\[
 \tau=\floor{E(m\,r_h+\delta)}+1,\qquad
 R=\operatorname{image}\bigl(\OO_K\to\OO_L/(\varpi_L^{\tau})\bigr),
\]
and let $V_j=((\xi^a-1)^j)_a\in(\OO_L/(\varpi_L^{\tau}))^{e_h}$, with $a$ ranging over the
units modulo $p^h$. Then $d_{K,h}(\delta)$ is the first $m$ with
\[
 -V_m\in R\,V_0+\cdots+R\,V_{m-1}.
\]
\end{proposition}

\begin{proof}
Since $E\,v$ takes integer values on $L^\times$, membership asserts precisely the existence
of a monic polynomial of degree $m$ over $\OO_K$ satisfying \eqref{cost:eq:monictest}.
\end{proof}

Only finite rings and linear membership occur. The coefficient ring is
$\OO_K/\mathfrak p_K^{\ceil{\tau/e(L/K)}}$, and it suffices to test one representative of each
$\Gal(L/K)$-orbit of test points.

The same test makes the aggregate recursion of Corollary~\ref{cost:cor:shells} exact over $K$.

\begin{proposition}[One polynomial over a prescribed field; \stU]\label{cost:prop:single}
Let $S$ and $\rho$ be as in Corollary~\ref{cost:cor:shells} and $\delta\ge0$. In increasing order of
$h\in S$ put
\[
 n_h=d_{K,h}\Bigl(\delta+\sum_{j\in S,\ j<h}n_j(r_j-r_h)\Bigr).
\]
Then the least value of $G_\rho(g)$ among nonzero $g\in\OO_K[T]$ that block
$\bigcup_{h\in S}Z_h$ at tolerance $\delta$ in the aggregate sense is $\sum_{h\in S}n_hr_h$.
\end{proposition}

\begin{proof}
Extract from a blocker $g$ its monic factor $P_h\in\OO_K[T]$ of roots of valuation $r_h$,
for each $h\in S$. By Lemma~\ref{cost:lem:roots},
$v(g(z))-G_\rho(g)=v(P_h(z))-(\deg P_h)r_h-(G_\rho(g)-G_{r_h}(g))$ on $Z_h$, and the savings
$G_\rho(g)-G_{r_h}(g)$ are at least $\sum_{j<h}(\deg P_j)(r_j-r_h)$, which by induction and
$r_j>r_h$ is at least the sum in the recursion. The extracted factor must therefore pass
the test \eqref{cost:eq:monictest} at that effective tolerance, so $\deg P_h\ge n_h$, and
\eqref{cost:eq:budget} gives the lower bound. Conversely the product of polynomials of
least degree for these tests has all roots on the designated radii, giving exactly those
savings and the claimed Gauss valuation.
\end{proof}

\subsection{Tame descent, cyclotomic and abelian coefficient fields}\label{cost:sec:abelian}
{\raggedright\sources{grh-torsion-and-prime-sampling; exponent_seven_and_grh_transfers; grh_sharp_torsion_and_symplectic_blocks}}

\begin{proposition}[Tame descent; \stU]\label{cost:prop:tame}
Let $K\subseteq L$ be finite extensions of $\Qp$ with $L/K$ tamely ramified, let $X$ be a
finite set of points stable under the absolute Galois group of $K$, and let $\theta$ be
real. If some monic $P\in\OO_L[T]$ of degree $m$ satisfies $v(P(z))>\theta$ for all $z\in X$,
then so does some monic polynomial of degree $m$ in $\OO_K[T]$. In particular
$c_{L,h}(\delta)=c_{K,h}(\delta)$ for all $h$ and $\delta$: every shell cost at every
tolerance is invariant under tame base change.
\end{proposition}

\begin{proof}
Choose $u\in\OO_L$ with $\Tr_{L/K}u=1$. Such an element exists: the trace of the maximal
unramified subextension is surjective on integers, because it reduces to the trace of
finite fields, and the remaining totally ramified degree is a unit; divide a trace-one
element of the unramified subextension by that degree. Take the trace of $uP$
coefficientwise. The result is integral and monic. Each conjugate of $P$ satisfies the
strict inequality at every point, by stability of $X$, and the ultrametric inequality gives
the same strict inequality for the trace. Extension of scalars gives the reverse
implication.
\end{proof}

The same argument works in several variables for a common evaluation threshold and an
imposed monic coefficient.

Fix a compatible system of primitive roots $\xi_h$ of order $p^h$ with $\xi_{h+1}^p=\xi_h$,
and let $K_s=\Qp(\xi_s)$, whose absolute ramification index is $e_s$. On an arbitrarily high
shell a packet over $K_s$ has the same optimum as a packet on shell $s$.

\begin{theorem}[Descent to the coefficient level; \stU]\label{cost:thm:cyclotomic}
Let $s\ge1$.
\begin{enumerate}[label=\textup{(\roman*)}]
\item $c_{K_s,h}(\delta)=c_{\min\set{h,s}}(\delta)$ for all $h\ge1$ and $\delta\ge0$, with
$c_j(\delta)$ as in \eqref{cost:eq:shellcost}. Hence
$C_{K_s}(S,\delta)=\sum_{h\in S}c_{\min\set{h,s}}(\delta)$, and for an interval starting at
$h_0\ge s$ the least endpoint guaranteeing a point of loss at most $\delta$ for every
packet of cost at most $B$ is $h_0+\floor{B/c_s(\delta)}$.
\item More precisely, let $h\ge s$ and let $P\in\OO_{K_s}[T]$ be monic of degree $m$ with all
roots of valuation $r_h$. Then $u=m\,e_s/e_h$ is an integer, and if $u<e_s$ some $z\in Z_h$
satisfies $v(P(z))-m\,r_h\le U_s(u)$.
\end{enumerate}
\end{theorem}

\begin{proof}
For $h\le s$ the field $K_s$ contains $\zeta_{p^h}$ and Theorem~\ref{cost:thm:packet} applies. Let
$h>s$. Partition the roots of $P$ into conjugacy orbits over $K_s$, repeated according to
multiplicity. If an orbit has degree $D$, its norm has valuation $D/e_h\in e_s^{-1}\Z$, so
$u_D=De_s/e_h$ is a positive integer. Pick a member $\alpha$ of the orbit and a nearest point
$\xi_h^b-1$ of the shell. Every automorphism fixing $K_s$ acts on $\xi_h$ by an exponent
congruent to $1$ modulo $p^s$, so each root of the orbit has a nearest point whose exponent
is congruent to $b$ modulo $p^s$; ties cause no problem, since one conjugating automorphism
can be chosen for each member. Assign $u_D$ marks to the unit class $b\bmod p^s$. If a
class $a\bmod p^s$ is different from $b$, with $j=v_p(a-b)<s$, then at every lift of $a$ to
the shell $h$ this orbit contributes to the excess at most
\[
 D\,\frac{p^j-1}{e_h}=u_D\,\frac{p^j-1}{e_s},
\]
by the nearest-point inequality of Lemma~\ref{cost:lem:shell}. The total number of marks is
$u=\sum u_D=m\,e_s/e_h$. If $u<e_s$, the least-occupied path on shell $s$ supplies an
unmarked class at which the total is at most $U_s(u)$, at every lift. This is (ii), and it
shows that passing the test \eqref{cost:eq:monictest} requires $u\ge m_s(\delta)$, that is,
$m\,r_h\ge c_s(\delta)$.

Conversely let $b$ run through the first $m_s(\delta)$ positive integers prime to $p$ and
take
\[
 P(T)=\prod_b\bigl((1+T)^{p^{h-s}}-\xi_s^b\bigr)\in\OO_{K_s}[T].
\]
The roots of the factor indexed by $b$ are the points $\xi_h^c-1$ with
$c\equiv b\pmod{p^s}$, and its value at $\xi_h^a-1$ is $\xi_s^a-\xi_s^b$. Hence the whole
excess profile of the extremal configuration at level $s$ pulls back to level $h$:
$v(P(z))-m\,r_h\ge U_s(m_s(\delta))>\delta$ on $Z_h$, with
$m\,r_h=m_s(\delta)\,r_s=c_s(\delta)$. The last assertion of (i) is the strict crossing of
the budget.
\end{proof}

For example, for $p=2$, $s=3$ and $\delta=1/4$ the exact price of every shell $h\ge3$ over
$\Q_2(\zeta_8)$ is $c_3(1/4)=1/2$, whereas a field containing all $2^h$-th roots of unity
permits the smaller costs of Example~\ref{cost:ex:packet} as $h$ grows. At $\delta=0$ the theorem
agrees with Theorem~\ref{cost:thm:field}, since $w=s-1$ and $c_j(0)=p^{1-j}$.

\begin{theorem}[All finite abelian coefficient fields; \stU]\label{cost:thm:abelian}
Let $K/\Qp$ be a finite abelian extension and $w=v_p(e_K)$. Write
$F_s=\Qp(\mu_{p^s})$.
\begin{enumerate}[label=\textup{(\alph*)}]
\item If $p$ is odd, then for all $h\ge1$ and $\delta\ge0$
\[
 c_{K,h}(\delta)=c_{\min\set{h,w+1}}(\delta),\qquad
 C_K(S,\delta)=\sum_{h\in S}c_{\min\set{h,w+1}}(\delta).
\]
In particular the entire tolerance profile depends only on the wild part of the
ramification index.
\item Let $p=2$. There are a finite unramified extension $U/\Q_2$ and a subfield $A$ of
some $F_n$ with $KU=AU$, and $c_{K,h}=c_{A,h}$ for all $h$. Let $s\ge1$ be minimal with
$A\subseteq F_s$. If $A=F_s$, then $c_{K,h}(\delta)=c_{\min\set{h,s}}(\delta)$. Otherwise
$s\ge3$, $[F_s:A]=2$, and
\[
 c_{K,h}(\delta)=
 \begin{cases}
 1,& h=1,\\[2pt]
 \dfrac1{e_h}\min\set{m\le e_h:E_h(m)>\delta},&2\le h<s,\\[8pt]
 \dfrac1{e_s}\min\set{m\in\set{0,2,4,\ldots,e_s}:U_s(m)>\delta},&h\ge s,
 \end{cases}
\]
where, for $2\le h<s$,
\[
 E_h(m)=\begin{cases}U_h(m),&m\text{ even},\\ U_h(m-1)+r_h-r_s,&m\text{ odd}.\end{cases}
\]
\end{enumerate}
\end{theorem}

\begin{proof}
By the local Kronecker--Weber theorem, $K\subseteq F_nU$ with $U/\Qp$ finite unramified
and $n\ge w+1$; see~\cite[Lecture~20]{Sutherland2017Number}. The two factors are linearly
disjoint, so $\Gal(F_nU/U)\simeq(\Z/p^n)^\times$. Moreover $K\cap U$ is the maximal
unramified subfield of $K$, so $[KU:U]=e_K$.

(a) For odd $p$ the group $(\Z/p^n)^\times$ is cyclic and $e_K$ divides $(p-1)p^w$.
Uniqueness of each intermediate degree in a cyclic extension gives
$KU\subseteq F_{w+1}U$. The extension $F_{w+1}U/KU$ has degree prime to $p$, and $KU/K$ and
$F_{w+1}U/F_{w+1}$ are unramified. By Proposition~\ref{cost:prop:tame} all the costs coincide with
those of $F_{w+1}$, which are given by Theorem~\ref{cost:thm:cyclotomic}.

(b) By the same compositum argument $KU=AU$ for the subfield $A\subseteq F_n$ fixed by the
corresponding subgroup, and Proposition~\ref{cost:prop:tame} gives $c_{K,h}=c_{A,h}$. If $A=F_s$,
Theorem~\ref{cost:thm:cyclotomic} applies. Assume $A\ne F_s$; then $s\ge3$. The group
$\Gamma=\Gal(F_s/A)$ avoids the nonidentity element $1+2^{s-1}$ of the kernel of reduction
from level $s$ to level $s-1$, by minimality of $s$. The unit group modulo $2^s$ is
$C_2\times C_{2^{s-2}}$; every element of order at least four has a power equal to
$1+2^{s-1}$, and among the four elements of order dividing two a subgroup avoiding
$1+2^{s-1}$ has order at most two. Hence $\Gamma$ has order two, generated by $\sigma=-1$ or
$\sigma=-1+2^{s-1}$ modulo $2^s$, and $e(A/\Q_2)=e_s/2$, so $v(A^\times)=2r_s\Z$.

\emph{Shells $h\ge s$.} By Proposition~\ref{cost:prop:monic}(i) the cost is $m/e_s$ with $m$ even, and
by Theorem~\ref{cost:thm:cyclotomic}(ii) a polynomial over $A\subseteq F_s$ of cost $m/e_s$ has
minimal excess at most $U_s(m)$. This is the lower bound. Conversely, for even $m\le e_s$
select the $m/2$ exponents $1+4i$, $0\le i<m/2$, and their images under multiplication by
$\sigma$. These are distinct, since $\sigma\equiv3\pmod4$. The marks are $\Gamma$-stable, and
each unit class modulo $2^j$ contains at least $\floor{m/e_j}$ of them: for $j\ge2$ the
first half cycles through all classes congruent to $1$ modulo $4$ with period $2^{j-2}$ and
the second half cycles through the complementary classes with the same period; at level one
every mark lies in the only class. The monic polynomial with roots $\xi_s^b-1$ at the
marks is therefore defined over $A$, and an unmarked point has excess at least $U_s(m)$, a
marked point infinite excess. The substitution $X\mapsto X^{2^{h-s}}$ in the variable
$X=1+T$ realizes the same cost at every level $h\ge s$.

\emph{Shells $2\le h<s$.} At level $h$ the automorphism $\sigma$ acts on exponents by
negation. For even $m$ the unrestricted upper bound for the minimal excess is $U_h(m)$, and
the configuration with $m/2$ marks in the class of $1$ modulo $4$ and their negatives
realizes it over $A$. For odd $m$, fix $t=\xi_s-1$ and $z=\xi_h-1$ and write $z=u+bt$ with
$u,b\in A$. Then
\[
 v(z-\sigma z)=2r_h,\qquad v(t-\sigma t)=2r_s,\qquad v(z-u)=v(bt)=2r_h-r_s=:\lambda,
\]
because $v_2(\sigma-1)=1$. The difference of two distinct points of $Z_h$ has valuation at
least $2r_h>\lambda$, so every point of $Z_h$ has distance exactly $\lambda$ from $u$. Since
$h<s$ gives $\lambda>r_h$, the single centre $u$ is integral, of valuation $r_h$, and
$T-u$ attains excess $r_h-r_s$ everywhere on the shell. Moreover $\lambda$ is an odd
multiple of $r_s$. If $[A(\alpha):A]$ is odd, the ramification index $e'$ of $A(\alpha)/A$
is odd and the value group $(2r_s/e')\Z$ does not contain $\lambda$. Consequently
$v(\alpha-u)\ne\lambda$ and $v(\alpha-z')\le\lambda$ for every $z'\in Z_h$. A monic
polynomial over $A$ of odd degree $m$ has an irreducible factor of odd degree; take one of
its roots $\alpha$ and apply the unrestricted bound $U_h(m-1)$ to the remaining roots. This
bounds its minimal excess by $E_h(m)$. Conversely multiply the extremal polynomial of even
degree $m-1$ by $T-u$; its minimal excess gains exactly $r_h-r_s$. All these polynomials
have every root on the designated radius. Finally the shell $h=1$ consists of the point
$-2$, and the positive integral degree is the cost.
\end{proof}

\begin{example}[Quadratic extensions of $\Q_2$]\label{cost:ex:quadratic}
For $K=\Q_2(\sqrt D)$ the theorem gives $c_{K,1}=1$ and, for $h\ge2$,
\[
 c_{K,h}(\delta)=
 \begin{cases}1/2,&0\le\delta<t_{D,h},\\1,&\delta\ge t_{D,h},\end{cases}
\qquad
\begin{array}{c|cc}
D&t_{D,2}&t_{D,h}\ (h\ge3)\\\hline
5&0&0\\
-1,\ 3&1/2&1/2\\
2,\ -2,\ 10,\ -10&1/4&1
\end{array}
\]
These radicands represent all seven quadratic extensions of $\Q_2$. Indeed
$\Q_2(\sqrt5)$ is unramified; $A=\Q_2(i)=F_2$ for $D=-1,3$, where
$U_2(1)=1/2$; and $A=\Q_2(\sqrt{\pm2})\subset F_3$ has index two for the remaining $D$,
where $E_2(1)=r_2-r_3=1/4$ and $U_3(2)=1$. Explicit extremizers in the variable $X=1+T$
are $X-i$ over $\Q_2(i)$ at $h=2$, the linear polynomial with centre $a=1+\sqrt{\pm2}$ at
$h=2$, for which $v(a^2+1)=3/2$, and
\[
 X^2-\sqrt2\,X+1,\qquad X^2-\sqrt{-2}\,X-1
\]
at $h=3$, whose root pairs are $\set{\xi,\xi^7}$ and $\set{\xi,\xi^3}$ for a suitable
primitive eighth root $\xi$; higher shells are obtained by $X\mapsto X^{2^{h-3}}$. For
instance, over $\Q_2(\sqrt{10})$ at $\delta=3/4$ the shell set $\set{2,3,4,5}$ costs
exactly $5/2$, while over $\Q_2(\sqrt3)$ the same four shells cost $4$; at $\delta=1$
every shell costs $1$ in every quadratic field. As an odd example, if $p=3$ and
$K/\Q_3$ is abelian with $w=1$, every shell $h\ge2$ costs $c_2(1/3)=2/3$ at tolerance
$1/3$, because $U_2(m)=\floor{m/2}/3$ for $m<6$; four consecutive shells therefore
guarantee a point of loss at most $1/3$ for a packet of cost at most $2$, and three can all
fail at exactly that cost. The dyadic constructions are checked by the delivered scripts
listed below; the first two are reported to validate $1{,}036$ balanced configurations,
exhausting the stable subsets through level five, and $84$ choices of degree, shell and
involution at conductors three through six.
\costfiles{
\item \artifact{exponent_seven_and_grh_transfers/other-results/check_balancing.py}
\item \artifact{exponent_seven_and_grh_transfers/other-results/check_lower_shells.py}
\item \artifact{exponent_seven_and_grh_transfers/other-results/eighth_roots.py}}
\end{example}

For primitive fourth roots the cost is known over every finite extension of $\Q_2$, abelian
or not. Define the truncated quadratic defect of $-1$ by
\[
 d_K=\max_{a\in\OO_K}\min\set{2,v(a^2+1)},\qquad\eta_K=\frac{d_K-1}2.
\]
The maximum exists since only the residue of $a$ modulo $4\OO_K$ matters, and $a=1$ gives
$d_K\ge1$.

\begin{proposition}[Primitive fourth roots; \stU]\label{cost:prop:fourth}
Let $p=2$ and let $K/\Q_2$ be finite. Then
\[
 c_{K,2}(\delta)=\begin{cases}1/2,&0\le\delta<\eta_K,\\ 1,&\delta\ge\eta_K.\end{cases}
\]
\end{proposition}

\begin{proof}
The two points of $Z_2$ are $i-1$ and $-i-1$. Write a linear centre as $a-1$. The
distances have valuations $v(a-i)$ and $v(a+i)$. Their sum is $v(a^2+1)$ and their minimum
is at most $1$, as the two targets differ by $2i$; if the minimum is less than $1$ the two
valuations are equal. Thus the minimum is exactly $\tfrac12\min\set{2,v(a^2+1)}$,
including $a=\pm i$. A monic linear polynomial passes the test \eqref{cost:eq:monictest}
precisely for $\delta<\eta_K$, and the monic quadratic with both targets as roots always
passes. These are all possible least degrees.
\end{proof}

For the quadratic fields the values are $d_K=1$ for $D=5$, $d_K=3/2$ for
$D=\pm2,\pm10$ and $d_K=2$ for $D=-1,3$, in agreement with Example~\ref{cost:ex:quadratic}; they
can be verified in sixteen residues per field
(\artifact{grh_sharp_torsion_and_symplectic_blocks/checks/results.txt}).

\begin{remark}[The cyclotomic intersection does not determine the cost]\label{cost:rem:twofields}
The fields $K_1=\Q_2(\sqrt3)$ and $K_2=\Q_2(\sqrt{10})$ have the same residue field and the
same value group, and both meet $\Q_2(\mu_{2^\infty})$ only in $\Q_2$: the three quadratic
subfields of the cyclotomic tower correspond to $-1,2,-2$, whose square classes differ from
those of $3$ and $10$. Nevertheless $c_{K_1,2}(1/4)=1/2$ and $c_{K_2,2}(1/4)=1$. In $K_1$
the centre $\sqrt3-1$ has distance $1$ from both targets. In $K_2$ the centre $\sqrt{10}$
has distance $3/4$ from both, since
$v((1+\sqrt{10})^2+1)=v(12+2\sqrt{10})=3/2$, and no element of $K_2$ improves this: its
difference from $\sqrt{10}$ has valuation in $\tfrac12\Z$ or is zero, and the ultrametric
inequality with unequal valuations bounds its distance to each target by $3/4$. Hence at
positive tolerance the shell cost is not a function of $K\cap\Qp(\mu_{p^\infty})$ and of
the admissible coefficient valuations.
\end{remark}

\subsection{Wild radical coefficient fields}\label{cost:sec:radical}
{\raggedright\sources{grh_transfer_extensions_2026-11-26; grh_method_transfers_graph_towers_and_finite_rings}}

The stabilization of the whole tolerance profile in Theorem~\ref{cost:thm:abelian} fails for
nonabelian wild extensions. Let $p$ be an odd prime, $u\in\Zp^\times$, and
\[
 K=\Qp(\pi),\qquad\pi^p=pu.
\]
The extension $K/\Qp$ is totally ramified of degree $p$, with ring of integers
$\Zp[\pi]$ and different exponent $2p-1$, since the derivative of $X^p-pu$ at $\pi$ is
$p\pi^{p-1}$. It is not Galois: normality would put $\zeta_p$ in a field of degree $p$.
For $h\ge2$ put
\[
 \theta_{p,h}=\frac1{p-1}-\frac{2p-1}{p\,e_h}>0 .
\]

\begin{theorem}[Best monic approximation; \stU]\label{cost:thm:approx}
Fix $h\ge2$, put $n=e_h/p$ and choose $\alpha\in Z_h$. For every $0\le m<e_h$,
\begin{equation}\label{cost:eq:approx}
 \max_{\substack{P\in\OO_K[T]\ \mathrm{monic}\\ \deg P=m}}v(P(\alpha))
 =\frac m{e_h}+\floor*{\frac mn}\theta_{p,h}.
\end{equation}
One monic polynomial $B$ of degree $n$ attains all these maxima: for $m=kn+j$ with
$0\le j<n$ and $0\le k<p$ use $T^jB(T)^k$. Every root of $B$ has valuation $r_h$, and
\[
 v(B(0))=\frac1p,\qquad
 v(B(\alpha))=\frac1p+\theta_{p,h}=\frac{(2p-1)(p^{h-1}-1)}{p\,e_h}.
\]
The valuations are the same at every $\alpha\in Z_h$. For $h=1$ the maximum in degree
$0\le m<e_1$ is $m/e_1$.
\end{theorem}

\begin{proof}
Write $C_j=\Qp(\zeta_{p^j})$ and $E_j=KC_j$; abbreviate $e=e_h$ and
$\theta=\theta_{p,h}$. Ramification breaks and differents are used as
in~\cite[Ch.~III--IV]{Serre1979Local}.

\emph{Ramification of the compositum.} The degree-$p$ extension $E_1/C_1$ is cyclic. The
coprime ramification indices $p$ and $p-1$ show that $E_1/\Qp$ is totally ramified of
degree $p(p-1)$. The element $(\zeta_p-1)/\pi$ has valuation $1/(p(p-1))$ and is a
uniformizer. A nonidentity element of $\Gal(E_1/C_1)$ multiplies it by a nontrivial $p$-th
root of unity, so the normalized valuation of the difference is $p+1$: the unique lower and
upper ramification break of $E_1/C_1$ is $p$. The extension $C_h/C_1$ has lower breaks
$p^j-1$, since for $v_p(a-1)=j$ the automorphism $\zeta\mapsto\zeta^a$ moves the
uniformizer $\zeta-1$ by an element of normalized valuation $p^j$; integrating the
reciprocal indices gives the upper breaks $p-1,2(p-1),\ldots,(h-1)(p-1)$. The extensions
$E_1/C_1$ and $C_h/C_1$ are Galois and linearly disjoint: their intersection, of degree
dividing $p$, could otherwise only be $E_1$, which would put $K$ inside the abelian
extension $C_h/\Qp$. Hence $[E_h:\Qp]=pe$ and $[E_h:K]=e$, and $\Gal(E_h/K)$ acts
transitively on $Z_h$. In the abelian group $\Gal(E_h/C_1)$ a character from the $E_1$
factor has break $p$ and the nontrivial characters from the other factor have breaks that
are multiples of $p-1$; these never coincide for odd $p$, and the product of two
characters with unequal breaks has the larger break. Consequently no nontrivial character
of the compositum is unramified, which proves total ramification and the product rule for
the upper filtration. The upper-to-lower map sends $p$ to
\[
 (p-1)+p\bigl(p-(p-1)\bigr)=2p-1,
\]
because the group index is $1$ up to $p-1$ and $p$ from $p-1$ to $p$. Lower ramification
groups intersect with subgroups, so the cyclic extension $E_h/C_h$ has its only lower break
at $2p-1$ and different exponent $2p(p-1)$.

\emph{The index of the power-basis order.} Let $d_C$ be the different exponent of
$C_h/\Qp$. The shifted cyclotomic polynomial is Eisenstein and
$\Zp[\alpha]=\OO_{C_h}$. Its degree remains $e$ over $K$, so the discriminant of the order
$\OO_K[\alpha]$ has valuation $p\,d_C$ in $K$. Transitivity of the different through $C_h$
and through $K$ gives the different exponent $2p(p-1)+p\,d_C-(2p-1)e$ of $E_h/K$. All
residue degrees are one, so the discriminant--index formula yields
\begin{equation}\label{cost:eq:index}
 I:=\operatorname{length}_{\OO_K}\bigl(\OO_{E_h}/\OO_K[\alpha]\bigr)
 =\frac{(2p-1)e-2p(p-1)}2.
\end{equation}

\emph{The best approximation in degree $n$.} Choose a monic $B\in\OO_K[T]$ of degree $n$
maximizing $b=v(B(\alpha))$. The maximum exists: the lower coefficients range over the
compact set $\OO_K^n$, and the image of the evaluation map is compact and does not contain
zero because $\alpha$ has degree $e>n$ over $K$. The polynomial $T^n$ gives $b\ge1/p$. We
claim $b\notin e^{-1}\Z$. Otherwise $b=j/e+t/p$ with $0\le j<n$ and $t\ge1$; both
$B(\alpha)$ and $\pi^t\alpha^j$ have this valuation, and since $E_h$ and $K$ have residue
field $\F_p$, a unit multiple of $\pi^tT^j$ cancels the initial residue of $B(\alpha)$ and
gives a better monic polynomial of the same degree.

The $e$ elements $w_{jk}=\alpha^jB(\alpha)^k$, $0\le j<n$, $0\le k<p$, form an
$\OO_K$-basis of $\OO_K[\alpha]$, by triangularity in their distinct degrees $kn+j$. Their
valuations are pairwise incongruent modulo $p^{-1}\Z$: the number $b$ has order $p$ in
$(pe)^{-1}\Z/e^{-1}\Z$, so equality of two classes forces equality of $k$ and then of $j$.
Every nonzero $K$-linear combination therefore has the valuation of its unique least term.
It follows that the elements $w_{jk}/\pi^{\floor{p\,v(w_{jk})}}$ form an integral basis of
$E_h/K$: their valuations lie in $[0,1/p)$ and are pairwise distinct, so a linear
combination is integral exactly when all its coefficients are. With $\theta'=b-1/p$, so
that $v(w_{jk})=(kn+j)/e+k\theta'$, the index is
\[
 I=\sum_{j,k}\floor{p\,v(w_{jk})}=\frac{(p-1)(e-1)}2+\frac{pe(p-1)}2\,\theta',
\]
because the fractional parts of the numbers $p\,v(w_{jk})$ are exactly
$0,1/e,\ldots,(e-1)/e$. Comparison with \eqref{cost:eq:index} gives $\theta'=\theta$.

\emph{All degrees.} In the monic triangular basis $T^jB^k$ every monic polynomial of degree
$m=kn+j$ has coefficient one at its leading basis element. The no-cancellation property
bounds its valuation at $\alpha$ by that of the leading basis element, which attains it.
This proves \eqref{cost:eq:approx}.

\emph{The roots of $B$.} Collect the roots of $B$ of valuation $r_h$ in a monic factor
$P\in\OO_K[T]$ of degree $a$. Every other root contributes at most $r_h$ to $v(B(\alpha))$.
If $a<n$, then \eqref{cost:eq:approx} in degree $a$ gives $v(P(\alpha))\le a/e$ and hence
$v(B(\alpha))\le n/e$, contrary to $\theta>0$. Therefore all roots of $B$ have valuation
$r_h$ and $v(B(0))=n/e=1/p$. The valuations are the same on the whole shell by the
transitivity of $\Gal(E_h/K)$.

\emph{The case $h=1$.} The values $v(a_j)+j/(p-1)$ of distinct nonzero coefficient terms of
degree below $p-1$ are incongruent modulo $p^{-1}\Z$, since $p$ and $p-1$ are coprime.
A monic polynomial of degree $m<p-1$ therefore cannot exceed $m/(p-1)$ in valuation, and
the monomial attains it.
\end{proof}

\begin{theorem}[Exact cost over wild radical fields; \stU]\label{cost:thm:radical}
Let $p$ be odd and $K=\Qp(\pi)$ with $\pi^p=pu$, $u\in\Zp^\times$. Then $c_{K,1}(\delta)=1$
and, for $h\ge2$ and $\delta\ge0$,
\[
 c_{K,h}(\delta)=\min\Bigl\{1,\ \frac{\floor{\delta/\theta_{p,h}}+1}p\Bigr\},
 \qquad C_K(S,\delta)=\sum_{h\in S}c_{K,h}(\delta).
\]
That is, a shell $h\ge2$ costs $k/p$ when
$(k-1)\theta_{p,h}\le\delta<k\theta_{p,h}$, $1\le k\le p-1$, and costs $1$ when
$\delta\ge(p-1)\theta_{p,h}$. The minimum is attained by the polynomials $B^k$ of
Theorem~\ref{cost:thm:approx} with witness $0$. The answer is independent of the unit $u$ and of
$\rho\ge\max_{h\in S}r_h$, and it is unchanged for bounded integral power series in
$\OO_K[[T]]$ in place of polynomials.
\end{theorem}

\begin{proof}
By Proposition~\ref{cost:prop:monic} it suffices to find $d_{K,h}(\delta)$. For $m<e_h$,
Theorem~\ref{cost:thm:approx} shows that the largest excess $v(P(\alpha))-m\,r_h$ of a monic
polynomial of degree $m$ is $\floor{m/n}\theta_{p,h}$, attained simultaneously at all
points of the shell. The least degree making it exceed $\delta$ is $kn$ with
$k=\floor{\delta/\theta_{p,h}}+1$, provided $k<p$, and it costs $kn\,r_h=k/p$. At degree
$e_h$ the shifted cyclotomic polynomial vanishes on the shell. The case $h=1$ follows from
the last part of Theorem~\ref{cost:thm:approx}. For bounded integral power series, truncate at a
degree so high that the omitted terms have valuation greater than $\sum_it_i+\delta$ at
every witness, at the radius $\rho$ and at every allowed torsion point; this is possible
because all these arguments have valuation bounded below by a positive number. The
thresholds are unchanged, each loss at most $\delta$ is unchanged and each loss exceeding
$\delta$ still exceeds it, so blocking is preserved and the polynomial bound applies.
\end{proof}

For fixed $p$ the thresholds $\theta_{p,h}$ increase strictly to $1/(p-1)$, so the
tolerance profile changes at every conductor, in contrast with
Theorem~\ref{cost:thm:abelian}(a); for each fixed $\delta$ the cost is eventually
$\min\set{1,(\floor{(p-1)\delta}+1)/p}$.

\begin{example}\label{cost:ex:radical}
For $p=3$ and $u=1$, that is $K=\Q_3(3^{1/3})$, one has
$\theta_{3,h}=\tfrac12-\tfrac5{3e_h}$, and on every shell $h\ge2$ the cost is $1/3$ for
$0\le\delta<\theta_{3,h}$, $2/3$ for $\theta_{3,h}\le\delta<2\theta_{3,h}$ and $1$ for
$\delta\ge2\theta_{3,h}$. At $h=2$ the polynomial $B(T)=T^2+\pi T+\pi$ is optimal:
\[
 \Res_Y\bigl(Y^3-3,\ \Res_T(\Phi_9(1+T),\,T^2+YT+Y)\bigr)=59049=3^{10},
\]
so $v(B(\zeta_9-1))=5/9=\tfrac13+\theta_{3,2}$, and the thresholds are $2/9$ and $4/9$. For
$p=5$, $h=2$ one has $\theta_{5,2}=4/25$ and the cost is $k/5$ on
$[(k-1)\tfrac4{25},k\tfrac4{25})$ for $k\le4$ and $1$ from $16/25$ on. Maximizing
polynomials constructed by successive residue cancellation for $u=1$, with exact
evaluation through the norm from $E_h$, give the following values, which use no
discriminant or ramification-break formula:
\[
\begin{array}{c|r|r|c|r}
(p,h)&e_h&\deg B&\theta_{p,h}&p\,e_h\,v(B(\alpha))\\\hline
(3,2)&6&2&2/9&10\\
(3,3)&18&6&11/27&40\\
(3,4)&54&18&38/81&130\\
(5,2)&20&4&4/25&36\\
(5,3)&100&20&29/125&216\\
(7,2)&42&6&6/49&78
\end{array}
\]
The certificates are the first two files below; the third records, for $p=3$ and
$2\le h\le6$, the scaled valuations $10,40,130,400,1210$.
\costfiles{
\item \artifact{grh_transfer_extensions_2026-11-26/data/false_tate_certificates.json}
\item \artifact{grh_transfer_extensions_2026-11-26/verification/false_tate_check.txt}
\item \artifact{grh_method_transfers_graph_towers_and_finite_rings/supplement/tower_results.txt}}
\end{example}

The fields $\Qp(\zeta_{p^h},p^{1/p})$ are the classical layers of the false Tate curve
extension; their ramification is computed in~\cite{Viviani2004Ramification}, explicit
uniformizers for $u=1$ are constructed in~\cite{WangYuan2023Uniformizer}, and the
maximizing monic polynomials are related to the reduced bases of the theory of higher
Newton polygons~\cite{GuardiaMontesNart2015Higher}.

\subsection{Simultaneous powers}\label{tors:sec:powers}
\sources{grh_beyond_dirichlet_grh_2026-11-25; grh_quantitative_twists_2026-11-25; grh_exact_torsion_transfer_2026-11-25; grh_sharp_transfers_2026-11-25; dirichlet_li_recovery}

The arithmetic construction needs one torsion tuple $\zeta$ that is good for the first
$m$ powers $\zeta,\zeta^2,\dots,\zeta^m$ of a character. All selection theorems transfer
to this situation by one substitution, at the price of replacing the depth $r$ by $r/m$
and the budget $B$ by $mB$.

\begin{lemma}[Power substitution]\label{tors:lem:powers}
Let $F\in\OO[[T_1,\dots,T_s]]$, let $x\in\Cp^s$ be a witness of depth $r>0$, and let
$j\ge1$. Put
\[
  F^{[j]}(T)=F\bigl((1+T_1)^j-1,\dots,(1+T_s)^j-1\bigr)\in\OO[[T_1,\dots,T_s]] .
\]
\begin{enumerate}[label=\textup{(\alph*)}]
\item There is $y\in\Cp^s$ with $v(y_a)\ge r/j$ for all $a$ and $F^{[j]}(y)=F(x)$. Thus
  $y$ is a witness of depth $r/j$ for $F^{[j]}$ with the same valuation, and
  $G_{r/j}(F^{[j]})\le v(F(x))$. If $F$ is a polynomial with coefficients in
  $\overline{\Z}_p$ and $x\in\overline{\Q}_p^{\,s}$, then $F^{[j]}$ has coefficients in
  $\overline{\Z}_p$ and $y$ can be chosen in $\overline{\Q}_p^{\,s}$.
\item For roots of unity $\zeta_a$ of $p$-power order,
  $F^{[j]}(\zeta_1-1,\dots,\zeta_s-1)=F(\zeta_1^j-1,\dots,\zeta_s^j-1)$.
\item If $\zeta$ has order $p^h$ and $p^h>m\ge j$, then $\zeta^j\ne1$ and
  $v(\zeta^j-1)=p^{\vp(j)}r_h\le m\,r_h$.
\end{enumerate}
\end{lemma}

\begin{proof}
(a) The substituted series have zero constant term, so $F^{[j]}$ is a well defined
integral series. For each coordinate choose $y_a$ with $(1+y_a)^j-1=x_a$ as follows. If
$x_a=0$ take $y_a=0$. Otherwise the polynomial $(1+Y)^j-1-x_a$ is monic and integral, so
its roots are integral, and their valuations sum to $v(x_a)\ge r$; at least one root has
valuation at least $r/j$; if $x_a$ is algebraic, so are these roots. Convergence of both
sides and the identity $F^{[j]}(y)=F(x)$ follow because all arguments have positive
valuation. (b) is the same identity at torsion
points. (c) follows from \eqref{tors:eq:shell}, since $p^{\vp(j)}\le j\le m<p^h$ and
$\zeta^j$ has order $p^{h-\vp(j)}$.
\end{proof}

\begin{proposition}[Selection for the first $m$ powers; \stU]\label{tors:prop:powers}
Let $(F_i)_{i\in I}$ be a finite family in $\OO[[T_1,\dots,T_s]]$ with witnesses $x_i$ of
depth $r>0$, let $t_i=v(F_i(x_i))$ with $\sum_it_i\le B$, and let $m\ge1$. In
\textup{(d)} let the $F_i$ moreover be polynomials with coefficients in $\overline{\Z}_p$
and $x_i\in\overline{\Q}_p^{\,s}$. Write $\zeta^j-1=(\zeta_1^j-1,\dots,\zeta_s^j-1)$.
\begin{enumerate}[label=\textup{(\alph*)}]
\item \emph{Nonvanishing.} With $H=h_*(r/m)+\floor{mB}$ there is a tuple with
  $\zeta_a^{p^H}=1$ and $0<v(\zeta_a-1)\le r/m$ such that $F_i(\zeta^j-1)\ne0$ for every
  $i$ and $1\le j\le m$.
\item \emph{Aggregate comparison.} Let $r_{h_0}\le r/m$. There is a tuple with
  $\zeta-1\in\torsZ{h_0}{H}^s$, $H=h_0+\floor{p^{h_0-1}mB}$, such that
  \[
    \sum_{i\in I}\sum_{j=1}^mv\bigl(F_i(\zeta^j-1)\bigr)\le mB .
  \]
\item \emph{Separate bounds, harmonic exponent.} Let $s\le s_0$, $\eps>0$,
  $\delta=\min\{\eps,1\}/(s_0+2)$, and let $H$ satisfy
  $mB+1<\delta\sum_{h=h_*(r/m)}^{H}(h+1)^{-1}$. There is a tuple with
  $\zeta-1\in\torsZ{h_*(r/m)}{H}^s$ such that
  \[
    v\bigl(F_i(\zeta^j-1)\bigr)\le t_i+\eps\qquad(i\in I,\ 1\le j\le m).
  \]
\item \emph{Fixed cardinality.} Let $t_i\le B_i$, $\eps_i>0$, $0<\eta\le1$,
  $s\le s_0$ and $r_{h_0}\le r/m$, and let $H$ satisfy \eqref{tors:eq:Hcondition} for the
  $m\cdot\#I$ constraints obtained by taking each budget $B_i+\eps_i$ and each tolerance
  $\eps_i/s_0$ with multiplicity $m$. There are at least $K^s$ tuples
  $\zeta-1\in\torsZ{h_0}{H}^s$, $K=\floor{(1-\eta)e_{h_0}}+1$, with
  $v(F_i(\zeta^j-1))\le t_i+\eps_i$ for all $i$ and $1\le j\le m$.
\end{enumerate}
If in \textup{(b)} or \textup{(d)} moreover $p^{h_0}>m$, then for every selected
coordinate and every $1\le j\le m$ one has $\zeta_a^j\ne1$ and $v(\zeta_a^j-1)\le r$.
\end{proposition}

\begin{proof}
By \cref{tors:lem:powers} the series $F_i^{[j]}$, $i\in I$, $1\le j\le m$, have
witnesses of depth $r/m$ with the valuations $t_i$, so their total budget is at most
$mB$, and their values at $\zeta-1$ are the numbers $F_i(\zeta^j-1)$. Apply
\cref{tors:thm:nonvanishing}, \cref{tors:cor:packets}(a), \cref{tors:thm:uniform} and
\cref{tors:thm:packet}, respectively, with depth $r/m$; in \textup{(c)} the total
positive excess at most $\eps$ given by \cref{tors:thm:uniform} bounds each excess. The
last assertion is \cref{tors:lem:powers}(c), since $m\,r_{h}\le m\,r_{h_0}\le r$ for
$h\ge h_0$.
\end{proof}

Part \textup{(c)}, for polynomials over $\overline{\Z}_p$ and witnesses with algebraic
coordinates, is the statement \code{exists_rootsOfUnity_ellVal_eval_le} used by the
source; the bound $H$ is chosen before the polynomials, and it is unchanged for series
over $\OO$ because of truncation and algebraic approximation. In \textup{(d)} the exponent is polynomial in the inverse
tolerances for fixed $\#I$, $m$ and $s_0$, and its degree grows with $m$. Part
\textup{(b)} is neither a prescribed-order statement nor a density statement for
arithmetic characters of bounded conductor.

\paragraph{Sharp endpoints for the first two powers.}
For the powers $J=\set{1,a,\dots,a^{k-1}}$ of an integer $a\ge2$ prime to $p$, put
$d_h=\ord_{p^h}(a)$. Every orbit of multiplication by $a$ on $\torsU_h$ is a cycle of length
$d_h$, and a cover in the sense of \cref{tors:thm:cover} must meet every cyclic run of $k$
consecutive positions. Its minimum size on a cycle is $\ceil{d_h/k}$: each marked position
covers at most $k$ starts, and the positions $0,k,2k,\dots<d_h$ work also across the seam.
Thus
\begin{equation}\label{tors:eq:geomcover}
  c_h=\frac{\ceil{d_h/k}}{d_h}.
\end{equation}
In particular, if $k\mid d_{h_0}$, then every window of $N$ consecutive shells beginning at
$h_0$ has exact cost $N/k$ in every dimension. The case $J=\set{1,2}$ settles the endpoint
of \cref{tors:prop:powers}(a) for $m=2$ at odd $p$.

\begin{theorem}[Sharp endpoints for the first two powers; \stU]\label{tors:thm:twopowers}
Let $p$ be odd, $J=\set{1,2}$, $h_0\ge1$ and $r\ge r_{h_0}$. Let $H_s(B)$ be the least
endpoint $H\ge h_0$ such that every admissible family \textup{(}as in
\cref{tors:thm:cover}\textup{)} of budget at most $B$ has a tuple
$z\in\torsZ{h_0}{H}^s$ with $F_i(z)\ne0$ and $F_i([2]z)\ne0$ for all $i$. Let
$d=\ord_p(2)$.
\begin{enumerate}[label=\textup{(\alph*)}]
\item If $d$ is even, $H_s(B)=h_0+\floor{2B}$ for every $B\ge0$.
\item If $d$ is odd, $H_s(k/2)=h_0+k-1$ for every integer $k\ge1$, and $H_s(B)=h_0$ for
  $0\le B<1/2$. For general $B$ with $\floor{2B}\ge1$,
  \[
    h_0+\floor{2B}-1\le H_s(B)\le h_0+\floor{2B}.
  \]
\item In one variable and for every $B$, $H_1(B)$ is the least $H$ such that
  $\sum_{h=h_0}^H\ceil{d_h/2}/d_h>B$.
\end{enumerate}
\end{theorem}

\begin{proof}
Write $\lambda=\vp(2^d-1)$. For odd $p$ the binomial theorem gives
$\vp((1+b)^p-1)=\vp(b)+1$ when $\vp(b)\ge1$, and exponentiation by an integer prime to $p$
leaves this valuation unchanged. It follows that $d_h=d\,p^{\max(0,h-\lambda)}$, so all the
orders $d_h$ have the parity of $d$. If $d$ is even, $c_h=1/2$ for all $h$ by
\eqref{tors:eq:geomcover}, every window of $N$ shells has exact cost $N/2$ by
\cref{tors:thm:cover}, and (a) follows.

If $d$ is odd, $c_h=1/2+1/(2d_h)$, so $Nc_H>N/2$ on every nonempty window. The possible
extra cost of the explicit blocker is uniformly bounded, since
\[
  \sum_{h\ge1}\frac1{d_h}=\frac{\lambda+1/(p-1)}{d}<1 ;
\]
indeed $p^\lambda\le2^d-1<p^d$ implies $\lambda\le d-1$. Thus a window of $N$ shells can be
blocked with a budget strictly less than $(N+1)/2$, and cannot be blocked with budget
$N/2$. For $B=k/2$ the window of $k$ shells always succeeds, whereas that of $k-1$ shells,
if nonempty, can be blocked. For arbitrary $B$ take $k=\floor{2B}$: a window of $k+1$
shells succeeds, while one of $k-1$ shells can be blocked. A single shell costs more than
$1/2$. Part (c) is the one-variable formula of \cref{tors:thm:cover}.
\end{proof}

\begin{remark}[The endpoint of \cref{tors:prop:powers}(a) for $m=2$]\label{tors:rem:twopowers}
In the normalization of \cref{tors:prop:powers}(a) take $h_0=h_*(r/2)$; then
$r_{h_0}\le r/2\le r$, and the sufficient endpoint $h_*(r/2)+\floor{2B}$ of that
proposition is $h_0+\floor{2B}$. By \cref{tors:thm:twopowers}, at a positive
half-integral budget this endpoint is optimal exactly when $\ord_p(2)$ is even, in every
dimension. Both cases occur for infinitely many primes: for $p\equiv7\pmod8$ the number $2$
is a square and $(p-1)/2$ is odd, so $\ord_p(2)$ is odd; for $p\equiv3\pmod8$ the number
$2$ is a nonsquare and $(p-1)/2$ is odd, so $\ord_p(2)$ is even. This follows from the
supplementary law of quadratic reciprocity and Dirichlet's theorem on primes in
progressions.
\end{remark}

\begin{example}\label{tors:ex:twopowers}
At $p=7$ one has $d_h=3\cdot7^{h-1}$, and the shell densities for $J=\set{1,2}$ are
$2/3$, $11/21$, $74/147$ for $h=1,2,3$. Take $h_0=1$. The first shell costs $2/3$ in every
dimension. On the first two shells the exact one-variable cost is $2/3+11/21=25/21$. In
every dimension a budget of $1$ already guarantees success on these two shells, since the
orbit bound gives $2\cdot11/21>1$, while the first shell alone can be blocked. At $p=5$
the first two shells together cost exactly $1$ in every dimension, and a third shell is
necessary at that budget. Exhaustive checks of $69$ geometric-progression covers and of
the unit orbit sizes of $512$ mixed tuples are recorded in
\artifact{dirichlet_li_recovery/checks/torsion.txt}.
\end{example}

\subsection{Matrices, determinants and Fitting ideals}\label{tors:sec:matrices}
\sources{grh_method_transfers_2026-11-25_1; grh_beyond_dirichlet_grh_2026-11-25; grh_quantitative_twists_2026-11-25; grh_exact_torsion_transfer_2026-11-25; grh_sharp_transfers_2026-11-25; grh_fixed_field_transfer_2026-11-25}

Every selection theorem applies to determinants. Let $A(T)$ be a $u\times v$ matrix with
entries in $\OO[[T_1,\dots,T_s]]$. Evaluation at a point of the open unit polydisc is a
ring homomorphism, so evaluating a formal minor agrees with taking the minor after
evaluation. For $1\le k\le\min\{u,v\}$ and such a point $x$ let
\[
  \mu_k(A(x))=\min\bigl\{v(D(x)):\ D\text{ a }k\times k\text{ minor of }A\bigr\}\in[0,+\infty].
\]
For the finitely presented module $\torscoker\bigl(A(x)\colon\OO^v\to\OO^u\bigr)$ the ideal
$\torsFitt_{u-k}$ is generated by the $k\times k$ minors, and finitely many elements of a
valuation ring generate the principal ideal of the member of least valuation; thus
$\mu_k(A(x))$ is the valuation of this Fitting ideal. No Noetherian hypothesis on $\OO$ is
involved.

\begin{proposition}[What a controlled minor gives]\label{tors:prop:matrices}
Let $D$ be a $k\times k$ minor of $A$, let $z\in\Cp^s$ have coordinates of positive
valuation, and suppose $\tau:=v(D(z))<\infty$.
\begin{enumerate}[label=\textup{(\alph*)}]
\item $\rank_{\Cp}A(z)\ge k$ and $\mu_k(A(z))\le\tau$. If $k$ is the rank of $A$ over the
  fraction field of $\OO[[T_1,\dots,T_s]]$, then $\rank_{\Cp}A(z)=k$.
\item If $A$ is square and $D=\det A$, then $A(z)$ is invertible and every entry of
  $A(z)^{-1}$ has valuation at least $-\tau$.
\item Suppose that the entries of $A(z)$ lie in $\OO_L$ for a finite extension $L/\Qp$
  with ramification index $e_L$, and that $k=u$. Then
  \[
    \frac{\torslen_{\OO_L}\torscoker A(z)}{e_L}\le\tau,
  \]
  with equality if $A$ is square.
\end{enumerate}
\end{proposition}

\begin{proof}
(a) The first two statements are immediate. Minors of size larger than the generic rank
are zero as formal series and vanish at every specialization. (b) The adjugate has
integral entries at $z$ and $A(z)^{-1}=\torsadj(A(z))/\det A(z)$. (c) For a square matrix
of full rank over the discrete valuation ring $\OO_L$, Smith normal form identifies the
length of the cokernel with the valuation of the determinant with respect to a
uniformizer of $L$, which is $e_L\,v$. For a rectangular matrix with $u$ rows, the
selected $u\times u$ submatrix has its image contained in the image of $A(z)$, so the
cokernel of $A(z)$ is a quotient of the cokernel of the submatrix.
\end{proof}

\begin{corollary}[Simultaneous specialization of selected minors; \stU]
\label{tors:cor:minors}
Let $(A_b)_{b\in I}$ be finitely many matrices over $\OO[[T_1,\dots,T_s]]$, of possibly
different sizes, and for each $b$ let $D_b$ be a nonzero $k_b\times k_b$ minor of $A_b$.
\begin{enumerate}[label=\textup{(\alph*)}]
\item \emph{Sum of denominators.} If $r_{h_0}\le\rho$ and $\sum_bG_\rho(D_b)\le B$, in
  particular if there are witnesses $x_b$ of depth $\rho$ with $\sum_bv(D_b(x_b))\le B$,
  there is $z\in\torsZ{h_0}{h_0+\floor{Bp^{h_0-1}}}^s$ with
  \[
    \sum_bv(D_b(z))\le\sum_bG_\rho(D_b)\le B .
  \]
  If the $A_b$ are square, $D_b=\det A_b$, and the specialized entries lie in $\OO_L$,
  then all $A_b(z)$ are nonsingular and
  $\sum_b\torslen_{\OO_L}\torscoker A_b(z)\le e_L\,B$.
\item \emph{Ranks.} If there are witnesses $x_b$ of depth $r$ with
  $\sum_bv(D_b(x_b))\le B$, then for every $m\ge1$ there is a tuple $\zeta$ with
  $\zeta_a^{p^H}=1$, $H=h_*(r/m)+\floor{mB}$, such that
  $\rank_{\Cp}A_b(\zeta^j-1)\ge k_b$ for all $b$ and $1\le j\le m$, with equality when
  $k_b$ is the generic rank of $A_b$.
\item \emph{Separate bounds.} Under the hypotheses of \cref{tors:thm:uniform} for the
  series $D_b$ and witnesses $x_b$, there is a tuple $z$ in the grid of that theorem with
  $\sum_b\bigl(v(D_b(z))-v(D_b(x_b))\bigr)_+\le\eps$; the index set may then be infinite.
  For square $A_b$ with $D_b=\det A_b$ all $A_b(z)$ are invertible and the entries of
  $A_b(z)^{-1}$ have valuation at least $-v(D_b(x_b))-\eps$. Under the hypotheses of
  \cref{tors:thm:packet} there are $K^s$ tuples with $v(D_b(z))\le v(D_b(x_b))+\eps_b$.
\item \emph{Separate exact bounds over a fixed field.} If all $A_b$ have entries in
  $\OO_K[[T_1,\dots,T_d]]$ for a finite extension $K/\Qp$ and $\sum_bG_\rho(D_b)\le B$,
  there is a tuple $z$ with exponents in $[h_0,H]$, $H$ as in \eqref{tors:eq:notower},
  such that $v(D_b(z))\le G_\rho(D_b)$ for every $b$.
\end{enumerate}
In each case \cref{tors:prop:matrices} applies to every $A_b$ at the selected tuple.
\end{corollary}

\begin{proof}
Apply \cref{tors:cor:packets}(a), \cref{tors:prop:powers}(a),
\cref{tors:thm:uniform,tors:thm:packet}, and \cref{tors:cor:notowerseries} to the series
$D_b$.
\end{proof}

The choice of a nonzero minor is an input. Identifying a particular Selmer module or
projection-denominator module with the cokernel of such a matrix, with integral
coefficients and a nonzero witness, is a separate arithmetic theorem which the selection
theorems do not supply.

\begin{corollary}[Fitting profiles; \stU]\label{tors:cor:fitting}
Let $(A_b)$ be finitely many matrices over $\OO[[T_1,\dots,T_s]]$, let $\mathcal K$ be a
finite set of pairs $(b,k)$, and for each $(b,k)\in\mathcal K$ let $x_{b,k}$ be a point of
depth $r>0$ with $\mu_k(A_b(x_{b,k}))\le\beta_{b,k}<\infty$. Let $m\ge1$, $\eps>0$, and
let $H$ be as in \cref{tors:prop:powers}\textup{(c)} with
$B=\sum_{(b,k)\in\mathcal K}\beta_{b,k}$. There is a single tuple $\zeta$ with
$\zeta-1\in\torsZ{h_*(r/m)}{H}^s$ such that
\[
  \mu_k\bigl(A_b(\zeta^j-1)\bigr)\le\mu_k\bigl(A_b(x_{b,k})\bigr)+\eps
  \qquad\bigl((b,k)\in\mathcal K,\ 1\le j\le m\bigr).
\]
The reference point may be different for each matrix and each size of minors.
\end{corollary}

\begin{proof}
For each $(b,k)$ choose a $k\times k$ minor $D_{b,k}$ attaining the minimum at
$x_{b,k}$ and apply \cref{tors:prop:powers}(c) to this family. At every point the minimum
over all $k\times k$ minors is at most the valuation of the chosen one.
\end{proof}

\begin{remark}\label{tors:rem:fitting}
(1) A finite matrix over a valuation ring can be diagonalized by repeatedly choosing a
nonzero entry of least valuation and clearing its row and column. If the nonzero diagonal
entries have valuations $\lambda_1\le\dots\le\lambda_l$, then
$\mu_k=\lambda_1+\dots+\lambda_k$. Thus \cref{tors:cor:fitting} controls cumulative
invariant-factor valuations. It does not bound each $\lambda_k$ separately: the lists
$(1,3)$ and $(0,4)$ have partial sums $1,4$ and $0,4$, while the last entry of the second
list is larger. Over the nondiscrete ring $\OO$ the number $\mu_k$ is not a module length.

(2) For a finite complex of finite free modules over $\OO[[T_1,\dots,T_s]]$, apply
\cref{tors:cor:minors}(b) to one minor of generic rank of each differential. Every
specialized complex over $\Cp$ then has the generic ranks, so its cohomology has the
generic dimensions; a generically exact complex becomes exact at each of the first $m$
powers of the selected tuple. This concerns cohomology over the field and says nothing
about integral torsion.
\end{remark}

\begin{example}\label{tors:ex:rank}
Let $p=3$, $r=1$, $m=3$, and
\[
  A(X,Y)=\begin{pmatrix}1&X\\ Y&XY+X\end{pmatrix},\qquad
  A'(X,Y)=\begin{pmatrix}X-Y&0\\ 0&1\end{pmatrix},\qquad F_3(X,Y)=(1+X)^3-1 ,
\]
so that $\det A=X$ and $\det A'=X-Y$. The witnesses $(3,3)$, $(3,6)$, $(3,3)$ give the
valuations $1,1,2$, so $B=4$, and \cref{tors:prop:powers}(a) guarantees the exponent
$h_*(1/3)+12=14$. The tuple $(\zeta_{27},\zeta_{27}^2)$, of exponent $3$, already works:
\begin{center}
\begin{tabular}{cccc}
\toprule
$j$ & $v(\det A)$ & $v(\det A')$ & $v(F_3)$\\
\midrule
1 & $1/18$ & $1/18$ & $1/6$\\
2 & $1/18$ & $1/18$ & $1/6$\\
3 & $1/6$ & $1/6$ & $1/2$\\
\bottomrule
\end{tabular}
\end{center}
All entries are finite, so both matrices have rank two at the three powers. The
valuations are verified with exact cyclotomic resultants in:
\torsfiles{
\item \artifact{grh_beyond_dirichlet_grh_2026-11-25/artifacts/exact_checks.json}}
\end{example}

\subsection{Packets of characters on one support}\label{tors:sec:characters}
\sources{grh_quantitative_twists_2026-11-25}

The positive fraction of good roots in \cref{tors:thm:fraction} gives many characters at
once in the arithmetic twisting construction of the source. We use the following input
from the source, stated in words; it is part of the construction of anticyclotomic twists
with small Selmer lengths described in \cref{src:sec:mechanism}, and it does not depend
on \cref{src:thm:dgrh}.

Fix the CM field $E$, the prime $p$, the set $S_0$ containing the places above $p$, an
inertia element $g_0$ at a place of $S_0$ away from $p$, the coefficient ring, the Selmer
data $T_i$ and the Stickelberger systems $Y_k$ of the source. Assume \code{IsTorsionLevel}
at the required torsion level and all fields of the structures \code{SelmerDatum} and
\code{StickelbergerSystem}, and write $\beta_a$ for their finite normalized valuation
budgets, including the global and local terms. In particular the Selmer datum has scalar
inertia of a special kind, local cohomological conditions and representation hypotheses;
these are not automatic for arbitrary elliptic curves (\cref{sec:sel}). The source
constructs an integer $s_0\ge1$, depending only on the fixed data, such that for all
rational $r>0$ with $s_0r<1$ and all sufficiently large $N$ there exist the following.
\begin{itemize}
\item An auxiliary set $Q$ of primes of degree one, disjoint from its conjugate
  $\overline Q$ and from the prescribed bad sets, with $2\#Q\le s_0$. It contains a
  designated first prime, so $Q$ is nonempty.
\item Characters $\lambda_u\colon G_E\to\Z/p^N\Z$ for $u\in Q$, unramified outside
  $u,\bar u$ and trivial above $p$, whose restrictions to inertia at $u$ are surjective.
  On inertia at $u$ every other $\lambda_{u'}$ is zero.
\item One polynomial $F_a$ with coefficients in $\overline{\Z}_p$ for each datum or
  system, in the variables $(T_u)_{u\in Q}$, and a witness of depth $r$ with algebraic
  coordinates at which its valuation is at most $\beta_a+s_0r$.
\end{itemize}
For $p^N$-th roots of unity $(\zeta_u)_{u\in Q}$ the character
$\psi=\prod_u\zeta_u^{\lambda_u}$ is anticyclotomic. The values of the $F_a$ at
$(\zeta_u-1)_u$ evaluate the relevant group-ring elements, or bound the normalized
twisted Selmer length when the valuation is strictly less than $\beta_a+1$; this is the
\code{TwistPolynomials} construction. The finishing lemma \code{isSmallTwist_twist} of the
source deduces the local and global conclusions (\code{IsSmallTwist}) from the polynomial
valuations and the bound for the coordinates.

\begin{theorem}[Bounded-order packets on one support; \stU]\label{tors:thm:arith}
Under the hypotheses just specified, fix $m\ge1$, $\eps>0$ and a number $N_0\ge1$ of
characters. There is an explicit exponent $H$, depending on $p,m,\eps,s_0,N_0$ and the
budgets $\beta_a$, with the following property. There are one nonempty auxiliary set $Q$
as above and at least $N_0$ distinct anticyclotomic characters $\psi$ such that:
\begin{enumerate}[label=\textup{(\roman*)}]
\item each $\psi$ has order dividing $p^H$;
\item every power $\psi^j$, $1\le j\le m$, is ramified at every prime of
  $Q\cup\overline Q$ and unramified outside it, is trivial at the places above $p$, and
  has the tame conductor $\prod_{u\in Q\cup\overline Q}u$;
\item every such power satisfies the \code{IsSmallTwist} inequalities of the source with
  tolerance $\eps$: normalized Selmer lengths at most $\beta_i+\eps$, Stickelberger
  valuations at most $\beta_k+\eps$, and sum of the valuations at the tame generators at
  most $\eps$.
\end{enumerate}
For fixed arithmetic data and fixed $m,N_0$ the exponent $H$ can be chosen polynomially
bounded in $1/\eps$ as $\eps\downarrow0$. No bound for the norms of the primes in $Q$ is
asserted.
\end{theorem}

\begin{proof}
Choose a rational $e$ with $\tfrac12\min\{\eps,\tfrac12\}\le e\le\min\{\eps,\tfrac12\}$
and put $r=e/(4s_0)$. We use \cref{tors:prop:powers}(d) with the budgets
$B_a=\beta_a+1$, the tolerances $\eps_a=e/2$, depth $r$, $\eta=\tfrac12$, and $s_0$ as the
bound for the number of variables. Choose $h_0$ so large that $r_{h_0}\le r/m$,
$p^{h_0}>m$ and $K=\floor{e_{h_0}/2}+1\ge N_0$, and choose $H$ by
\eqref{tors:eq:Hcondition} for the constraints of \cref{tors:prop:powers}(d). All these
choices precede the polynomials. If there are no constraints the selection is trivial.

Next choose $N\ge H$ so large that the construction of the source exists. The
polynomials $F_a$ have witnesses of depth $r$ with valuations
$t_a\le\beta_a+s_0r<\beta_a+1$. \Cref{tors:prop:powers}(d) gives at least
$K^{\#Q}\ge N_0$ tuples $(\zeta_u)_{u\in Q}$ such that every power $j\le m$ has
$\zeta_u^j\ne1$, $v(\zeta_u^j-1)\le r$, and
\[
  v\bigl(F_a((\zeta_u^j-1)_u)\bigr)\le t_a+\tfrac e2\le\beta_a+\tfrac{3e}4<\beta_a+1 .
\]
The finishing lemma of the source applies with tolerance $e$, since $s_0r=e/4\le e<1$,
and monotonicity raises the tolerance to $\eps$. Every $\zeta_u$ is a $p^H$-th root of
unity, hence a $p^N$-th root of unity, so $\psi^{p^H}=1$ even if the auxiliary level $N$
is much larger than $H$.

For each $u\in Q$ fix an element $g_u$ of inertia above $u$ with $\lambda_u(g_u)=1$ and
$\lambda_{u'}(g_u)=0$ for $u'\ne u$; it does not depend on the tuple. Then
$\psi(g_u)=\zeta_u$, so different tuples give different characters, and the order of
$\psi$ is the largest order of a coordinate. Each power $\psi^j$ is nontrivial on
inertia at each $u$, and since $\psi$ is anticyclotomic its value on the conjugate
inertia is the inverse, so $\psi^j$ is ramified at $\bar u$ as well. Outside these primes
it is unramified. The primes are prime to $p$ and the characters have $p$-power order,
so the ramification is tame, and a nontrivial tame character has conductor exponent one.
This proves the statement about conductors. The polynomial dependence on $1/\eps$ follows
because the number of polynomials is fixed and $h_0$ grows only logarithmically in $1/e$
for fixed $N_0$ and $m$.
\end{proof}

\begin{remark}\label{tors:rem:arith}
(1) The characters of \cref{tors:thm:arith} have a common support, chosen after the
size $N_0$ of the packet is specified, one common auxiliary level $N$ and one
parametrization. The theorem adds an explicit bound for the order and a multiplicity to
the qualitative existence theorem of the source. It is not a statement of positive
density when conductors are ordered by norm, and a bound for the order of a character
does not bound the norm of its conductor. The orders within the packet may vary.

(2) Small normalized length is not vanishing. If the coefficient field has absolute
ramification index $e_L$, a nonzero module of finite length has normalized length at
least $1/e_L$, and $e_L$ grows with $H$. Producing for every $\eps>0$ a twist with
normalized Selmer length at most $\eps$ does not produce a zero Selmer module; the
resulting ramification obstruction is discussed in \cref{sec:sel}. Improving $H$ does not
remove it.
\end{remark}

\begin{corollary}[Characters of prime conductor; \stG]\label{cost:cor:cardchar}
Under the hypotheses of Theorem~\ref{cost:thm:cardmulti} write $H=H_n(B)$. There are a
prime
\[
 q\equiv1\pmod{p^H},\qquad q<2\bigl(p^HH\log p\bigr)^2,
\]
a generator $u$ of $(\Z/q\Z)^\times$, and nonprincipal Dirichlet characters
$\chi_1,\dots,\chi_d$ modulo $q$ with values in $\Cp$, of $p$-power orders between
$p^{h_0}$ and $p^H$ and of conductor $q$, such that
\[
 v\bigl(F_i(\chi_1(u)-1,\dots,\chi_d(u)-1)\bigr)\le A_i\qquad(1\le i\le n),
\]
and likewise with separate witnesses as in Theorem~\ref{cost:thm:cardmulti}.
\end{corollary}

\begin{proof}
Under the Riemann hypothesis for Dirichlet $L$-functions, supplied by \cref{src:thm:dgrh},
Bach and Sorenson's bound \cite{Bach1996Explicit} gives a prime $q\equiv1$ modulo $p^H$
below $2(\varphi(p^H)\log p^H)^2\le2(p^HH\log p)^2$. The group $(\Z/q\Z)^\times$ is cyclic,
and $p^H\mid q-1$, so $\chi_j(u)=1+z_j$, with $z_j$ from Theorem~\ref{cost:thm:cardmulti},
defines a character of order between $p^{h_0}$ and $p^H$. It is nonprincipal because
$h_0\ge1$, so its conductor is the prime $q$.
\end{proof}

The conductor bound is not claimed to be sharp, whereas the local exponent is. The
corollary realizes the zero-tolerance specialization by Dirichlet characters but imposes no
further Frobenius, ramification or Selmer conditions.

\subsection{When local torsion data become global twists}\label{spn:sec:characters}
\sources{grh_transfer_bridges_2026-11-25}

A torsion-avoidance lemma (\cref{sec:tors}) chooses roots of unity satisfying
polynomial inequalities. A global twisting theorem needs those roots to be
values of an arithmetic character. There are three separate quantities: the
orders of the local inertia values, the order of the global character, and
its conductor. The following elementary theorem identifies their first
compatibility obstruction.

\begin{theorem}[Exact extension criterion; \stU]\label{spn:thm:extension}
Let $N\ge1$, and let
\[
 0\longrightarrow A\longrightarrow B\longrightarrow C\longrightarrow0
\]
be an exact sequence of finite abelian $p$-groups, with $p^aC=0$. Let $D$ be
a subgroup of $B$, and let $\lambda\colon A\to\mu_{p^N}$ be a character.
\begin{enumerate}
\item An extension $\chi\colon B\to\mu_{p^{N+a}}$ trivial on $D$ exists if
and only if
\begin{equation}\label{spn:eq:compat-raised}
 \lambda(A\cap D)=1 .
\end{equation}
\item An extension $\chi\colon B\to\mu_{p^N}$ trivial on $D$ exists if and
only if
\begin{equation}\label{spn:eq:compat-exact}
 \lambda\bigl(A\cap(D+p^NB)\bigr)=1 .
\end{equation}
\end{enumerate}
Here $\lambda(H)=1$ means that $\lambda$ is trivial on every element of $H$.
\end{theorem}

\begin{proof}
A character from a subgroup of a finite abelian $p$-group extends to
$\mu_{p^\infty}$. An elementary construction adjoins one generator $b$ at a
time: if the least positive integer $r$ with $rb$ in the old subgroup is
$r=p^j$, choose an $r$th root of the already assigned value at $rb$. This
gives a well-defined extension and terminates.

Condition \eqref{spn:eq:compat-raised} is necessary for an extension killing
$D$. If it holds, $\lambda$ descends to $\bar A=(A+D)/D\subseteq B/D$. Extend
it into $\mu_{p^\infty}$. The quotient $(B/D)/\bar A$ is a quotient of $C$,
hence is killed by $p^a$. For each $\bar b\in B/D$,
\[
 \chi(\bar b)^{p^{N+a}}=\lambda(p^a\bar b)^{p^N}=1 .
\]
Thus the extension takes values in the required finite group.

For the second assertion, a character of exponent dividing $p^N$ kills $p^NB$
as well as $D$, giving necessity. Conversely, under
\eqref{spn:eq:compat-exact}, $\lambda$ descends to the image of $A$ in
$B/(D+p^NB)$. Extend it as above. This quotient is killed by $p^N$, so its
character values automatically lie in $\mu_{p^N}$.
\end{proof}

The permitted increase in order is bounded by the exponent of the quotient,
not by its cardinality. When $D=0$, the proportion of local
$\mu_{p^N}$-valued characters that extend with that same exponent is exactly
\begin{equation}\label{spn:eq:proportion}
 \frac{1}{[A\cap p^NB:p^NA]} .
\end{equation}
Indeed, all such local characters form the dual of $A/p^NA$, while the
extendible characters form the dual of $A/(A\cap p^NB)$. An argument that
samples the unrestricted local torsion grid cannot assume that the same
sampling measure applies to global characters.

Let $E$ be a CM field, with complex conjugation $c$. Fix an odd prime $p$ and
a squarefree modulus
\[
 \mathfrak m=\prod_{j=1}^t\mathfrak v_j\,c(\mathfrak v_j),
\]
where all the primes are distinct, $\mathfrak v_j\ne c(\mathfrak v_j)$, and
none lies above $p$. Put $q_j=\Nm\mathfrak v_j$. For a finite abelian
$p$-group with $c$-action, the minus part is the image of the idempotent
$(1-c)/2$.

\begin{theorem}[CM ray realization; \stU]\label{spn:thm:ray}
Set
\[
 R=\prod_{j=1}^t\F_{q_j}^{\times}[p^\infty],\qquad
 B=\Cl_{\mathfrak m}(E)[p^\infty]^{-},\qquad C=\Cl(E)[p^\infty]^{-}.
\]
Let $W$ be the image of $\mu(E)[p^\infty]$ in the minus local residue group
$R$. Then $W$ is cyclic, and there is an exact sequence
\begin{equation}\label{spn:eq:ray}
 0\longrightarrow R/W\longrightarrow B\longrightarrow C\longrightarrow0 .
\end{equation}
Consequently \cref{spn:thm:extension} applies with $A=R/W$. A proposed local
tuple must first annihilate $W$. Prescribed splitting conditions are imposed
by the appropriate subgroup $D$ of ray classes.
\end{theorem}

\begin{proof}
Start with the usual ray class exact sequence
\[
 \OO_E^\times\longrightarrow(\OO_E/\mathfrak m)^\times\longrightarrow
 \Cl_{\mathfrak m}(E)\longrightarrow\Cl(E)\longrightarrow0 .
\]
The Chinese remainder theorem identifies each conjugate residue pair with two
copies of $\F_{q_j}^\times$, exchanged by $c$. Its minus $p$-part consists of
the pairs $(u,u^{-1})$ and is one copy of $\F_{q_j}^\times[p^\infty]$.
Primary decomposition and the minus idempotent are exact direct-summand
operations on the finite groups.

The minus part of the unit image needs care. For any unit $u$, the quotient
$u/c(u)$ is an algebraic integer unit all of whose conjugates have absolute
value one. Kronecker's theorem makes it a root of unity. On the $p$-part,
taking its square root is an automorphism, since $p$ is odd. Thus the minus
unit image is contained in the image of $\mu(E)[p^\infty]$. The reverse
inclusion holds because every such root is a unit and $c$ acts on it by
inversion. The unit image is therefore exactly $W$, proving
\eqref{spn:eq:ray}.
\end{proof}

The ideal-theoretic class-field correspondence \cite{Milne2020CFT} interprets
characters of $B$ as finite anticyclotomic Hecke characters. If
$p\nmid\abs{\mu(E)}$, then $W=0$ and the local variables are unrestricted
before the class and splitting conditions are imposed. If $p\nmid h_E$, then
additionally $C=0$, so the unrestricted minus local tuple determines a unique
$p$-primary ray character. Neither simplification can be assumed in general.
In particular, the field of the source contains a primitive root of unity of
order $5040$; its cases $p=3,5,7$ require the $W$ term.

An extension supplied by \cref{spn:thm:ray} has conductor dividing
$\mathfrak m$. At each prime of this squarefree modulus, a nontrivial local
residue character has conductor exponent one. Therefore, if all prescribed
local characters are nontrivial,
\begin{equation}\label{spn:eq:conductor}
 \mathfrak f(\chi)=\mathfrak m,\qquad
 \Nm\mathfrak f(\chi)=\prod_{j=1}^tq_j^2 .
\end{equation}
Increasing the allowed global order does not force an increase in this
conductor. If the local character at pair $j$ has exact order $p^{n_j}$, the
first $M$ positive powers of $\chi$ all retain the full conductor if and only
if $p^{n_j}>M$ for every $j$. This assertion concerns local orders, even when
the global order is larger.

\begin{proposition}[A fixed-exponent parameter bound; \stU]\label{spn:prop:grid}
Suppose $N\ge1$ and $R=\prod_{j=1}^tC_{p^{h_j}}$ with $h_j\ge N$. Let $C$
require $r$ generators, let $W$ require $w$ generators, and let the
required-triviality subgroup $D\subseteq B$ have a chosen generating set of
size $s$. Then the image of
\[
 \Hom(B/D,\mu_{p^N})\longrightarrow\Hom(R,\mu_{p^N})
\]
contains a free $\Z/p^N\Z$-submodule of rank at least $\max(0,t-w-r-s)$. For
the CM unit subgroup, $w$ is zero or one.
\end{proposition}

\begin{proof}
Choose invariant-factor generators $c_i$ of $C$, of orders $p^{a_i}$, and
lifts $b_i$ in $B$. Write each relation $p^{a_i}b_i$ as an element of $R/W$,
and choose a representative in $R$. Assign the value zero to every $b_i$,
writing characters additively in $\Z/p^N\Z$. Assigning values to the $t$
local generators then gives a character precisely when the unit relations and
these $r$ relations are respected. Requiring the chosen $s$ elements of $D$
to vanish adds at most $s$ further linear equations. Thus the kernel of a
matrix with at most $w+r+s$ rows over $\Z/p^N\Z$ maps into the restriction
image. Smith normal form leaves at least the asserted number of free
columns.
\end{proof}

The bound $t-r$ in the case $W=D=0$ is sharp. For $a\ge1$, take
\[
 B=C_{p^{N+a}}^{\,r}\times C_{p^N}^{\,t-r},\qquad
 A=p^aC_{p^{N+a}}^{\,r}\times C_{p^N}^{\,t-r}.
\]
The first $r$ restriction coordinates lie in $p^a(\Z/p^N\Z)$ and have smaller
exponent; the last $t-r$ give all free coordinates of exponent $p^N$. The
proposition does not say that every original inertia coordinate can vary
independently, or even that every one is nonzero somewhere on the displayed
submodule. A polynomial which is nonzero on the full local parameter space
may vanish identically on the compatible subgroup; applying a
torsion-avoidance result after this restriction requires a fresh witness
estimate.

%======================================================================

\subsection{An explicit conductor bound and its limits}\label{spn:sec:conductor}
\sources{grh_transfer_bridges_2026-11-25}

The source's torsion and prime-selection arguments control the number and
character-theoretic properties of auxiliary primes. They do not give a norm
bound for primes satisfying every required Frobenius and Selmer condition.
For the simpler, unrestricted realization problem there is a completely
explicit bound.

\begin{theorem}[Bounded conductor for an abelian CM base; \stG]\label{spn:thm:bounded-conductor}
Let $E/\Q$ be an abelian CM extension contained in $\Q(\zeta_f)$. Let $p$ be
odd, $N\ge1$, and $t\ge1$. Write
\[
 \abs{\mu(E)[p^\infty]}=p^{w_0},\qquad p^a\,\Cl(E)[p^\infty]^{-}=0 .
\]
Let $S$ be a finite set of rational primes to avoid. Set
\[
 M=\lcm\Bigl(f,\ p^{N+w_0},\ \prod_{q\in S}q\Bigr).
\]
Choose a prime $\ell>t$ not dividing $M$, and put
\[
 Q=M\ell,\qquad P=(\varphi(Q)\log Q)^2 .
\]
Then, by Dirichlet GRH (\cref{src:thm:dgrh}), there exist $t$ distinct
rational primes $q_j\le P$, splitting completely in $E$ and avoiding $S\cup\set{p}$, with the following property.
Choose one conjugate pair above each $q_j$. Every tuple of local characters
of order dividing $p^N$ on those pairs extends to an anticyclotomic character
$\chi$ with
\[
 \ord(\chi)\mid p^{N+a},\qquad\Nm\mathfrak f(\chi)\le P^{2t}.
\]
If all the selected local characters have exact order $p^N$, then the
conductor is exactly the selected modulus, and the first $m$ powers of $\chi$
have that same conductor whenever $m<p^N$.
\end{theorem}

\begin{proof}
Choose $t$ distinct nonzero residues modulo $\ell$. For each, the Chinese
remainder theorem gives a reduced residue class modulo $Q$ which is $1$
modulo $M$ and has the chosen value modulo $\ell$. Under Dirichlet GRH, the
least-prime bound $P(a,q)\le(\varphi(q)\log q)^2$ for $q>3$ of Lamzouri, Li
and Soundararajan \cite[Corollary~1.2]{LamzouriLiSoundararajan2015Conditional} gives a prime
in each class at most $P$. They are distinct because the residues modulo $\ell$ are distinct.
Their congruence modulo $f$ ensures complete splitting in $E$; their
congruence modulo the other factors ensures avoidance and
$\vp(q_j-1)\ge N+w_0$.

For the corresponding local $p$-group $R$, every coordinate of the root-unit
image $W$ has order at most $p^{w_0}$. In a cyclic factor of order
$p^{h_j}$, with $h_j\ge N+w_0$, such elements lie in
$p^{h_j-w_0}R_j\subseteq p^NR_j$. Therefore $W\subseteq p^NR$. Every local
character into $\mu_{p^N}$ automatically annihilates $W$. Apply
\cref{spn:thm:extension,spn:thm:ray} with $D=0$. The conductor and power
assertions follow from \eqref{spn:eq:conductor}.
\end{proof}

This strengthens a construction which simply excludes primes dividing the
order of the roots of unity of $E$: deeper local residue groups absorb that
unit relation. Among supports with $h_j\ge N$, the extra depth
$h_j\ge N+w_0$ is also necessary for \emph{every independent} local
$\mu_{p^N}$ tuple to kill the unit image. Reduction away from $p$ injects
$\mu(E)[p^\infty]$ into each residue field, so the projection of $W$ is
exactly the subgroup of order $p^{w_0}$; it lies in $p^NR_j$ precisely at the
displayed depth, and a tuple supported on one coordinate detects any failure.
The construction also supplies arbitrarily many independent local coordinates
using distinct rational primes, without an exponentially growing recursive
exclusion modulus.

There is an elementary lower restriction. A tame inertia character of order
$p^N$ requires $p^N\mid q_j-1$, hence $q_j\ge p^N+1$. A nontrivial conjugate
pair therefore has conductor norm at least $(p^N+1)^2$. No conductor bound
independent of the requested local order can hold.

\begin{lemma}[The tame part of the small-valuation condition; \stU]\label{spn:lem:tame-cost}
Normalize $\vp(p)=1$. Let $\chi$ be a character as in \cref{spn:thm:ray}
whose local values on all $t$ conjugate pairs of $\mathfrak m$ have exact
order $p^N$, and let $\tau_{\mathfrak v}$ be a tame inertia generator at
$\mathfrak v\mid\mathfrak m$. Then
\[
 \sum_{\mathfrak v\mid\mathfrak m}\vp\bigl(1-\chi(\tau_{\mathfrak v})^j\bigr)
 \le\frac{2t\,p^{\floor{\log_pm}}}{(p-1)p^{N-1}},\qquad1\le j\le m<p^N .
\]
\end{lemma}

\begin{proof}
At a primitive $p^N$th root, $\vp(1-\zeta_{p^N})=1/((p-1)p^{N-1})$, and for
$1\le j\le m<p^N$,
$\vp(1-\zeta_{p^N}^{\,j})=p^{\vp(j)}/((p-1)p^{N-1})$. Primitivity makes the
valuation independent of the choice of generator. Sum over the $2t$ primes.
\end{proof}

Choosing $N$ large makes this particular cost small, with the conductor bound
of \cref{spn:thm:bounded-conductor}; this recovers one quantitative component
of the source's twisting problem. The theorem sets $D=0$: it does not
impose splitting at the primes above $p$, and \cref{spn:ex:conductor13}
below shows that this omission can be decisive even at tiny discriminant. Nor
do the congruences on $q_j$ enforce the Frobenius neighborhoods governing the
source's tied Selmer relations, which can involve nonabelian extensions. One
may quotient by the required subgroup $D$ and apply
\cref{spn:thm:extension}, or use \cref{spn:prop:grid} to exhibit a surviving
family; either route still needs a proof that the determinant or
special-value witness does not vanish on that family. Prescribed tame
ramification through relations among Frobenius elements is a developed
subject, including the Gras--Munnier criterion \cite{Hajir2022Tame}, and
effective Chebotarev results \cite{Thorner2017Explicit} can be considered
once the governing field, its degree and its discriminant are bounded; they
cannot be replaced by the Dirichlet endpoint \cref{src:thm:dgrh}. For a
nonabelian base field, even finding small completely split primes requires
additional analysis, and for $p=2$ the minus idempotent used here is
unavailable.

%======================================================================

\subsection{Exact conductor minima over \texorpdfstring{$\Q(\sqrt{-23})$}{Q(sqrt(-23))} and the precursor condition}\label{spn:sec:benchmark}
\sources{grh_transfer_bridges_2026-11-25; grh_ideal_densities_and_transfer_conditions_2026-11-26}

The source's predicate \code{IsFromU1} asks for an auxiliary character of the
form
\[
 \psi=\xi/\xi^c ,
\]
where the \emph{precursor} $\xi$ is ramified on only one side of each
conjugate pair and is itself trivial on the whole decomposition groups, not
only the inertia subgroups, at the prescribed places
(\code{IsFromU1.eq_one_of_mem_decompGroup}). Triviality of the output $\psi$
alone does not imply this stronger assertion. This subsection makes the
distinction effective in a small benchmark, computes the exact conductor
minima under both conditions, and realizes the precursor condition with
arbitrarily small normalized tame cost.

Let
\[
 K=\Q(\sqrt{-23}),\qquad\theta=\frac{1+\sqrt{-23}}2,\qquad
 \theta^2-\theta+6=0,\qquad\beta=1+2\theta=2+\sqrt{-23},
\]
and let $c$ be complex conjugation. The reduced primitive forms of
discriminant $-23$ are $(1,1,6)$, $(2,-1,3)$, $(2,1,3)$; the reduced-form
inequality $a\le\sqrt{23/3}<3$ makes this list exhaustive. Hence
$\Cl(K)\simeq C_3$, entirely in the minus part, and the only roots of unity
are $\pm1$. The prime $\mathfrak s=(3,\theta-1)$ is nonprincipal: the norm
form is $u^2+uv+6v^2$; for $v\ne0$ it is at least $23v^2/4>3$, and for $v=0$
it is a square. Since $\Nm\beta=27$ and $\beta$ reduces to zero at
$\mathfrak s$ and to a unit at the conjugate prime above $3$,
\begin{equation}\label{spn:eq:class-relations}
 \mathfrak s^3=(\beta),\qquad\mathfrak s\,c\mathfrak s=(3).
\end{equation}
In particular the class of $\mathfrak s$ generates $\Cl(K)$. Let
$S=\set{\mathfrak s,c\mathfrak s}$. The $S$-units are exactly
$\set{\pm3^a\beta^b:a,b\in\Z}$: their valuation vectors $(u,v)$ satisfy
$u-v\equiv0\pmod3$, and $(1,1),(3,0)$ generate precisely that lattice.

\begin{example}[Conductor $13$: a necessary increase of order, and a fatal splitting condition]\label{spn:ex:conductor13}
At $13$, the roots of $T^2-T+6$ are $5$ and $9$, so $13$ splits. The minus
local $3$-group has order three, and the ray exact sequence \eqref{spn:eq:ray}
gives a group $B$ of order nine,
$0\to C_3\to B\to C_3\to0$. At the root $\theta=5$, the residue of
$\beta/c(\beta)$ is $4$ modulo $13$, and $4^4\equiv9\not\equiv1\pmod{13}$.
Thus its $3$-primary component is nontrivial. By
\eqref{spn:eq:class-relations}, the triple of a minus lift of the class above
$3$ is a nonzero local inertia element. The lift therefore has order nine and
$B\simeq C_9$.

Its local inertia subgroup is $3C_9$. Every character $C_9\to\mu_3$ kills
that subgroup. A prescribed nontrivial local cubic character therefore needs
global order nine, even though its conductor norm remains $13^2=169$: the
order increase in \cref{spn:thm:extension} can be necessary in an actual
number field. The splitting condition above $3$ is more restrictive. The
minus Frobenius at $3$ generates $C_9$; requiring its value to be one kills
the whole ray group. Thus no nontrivial character of any $3$-power order on
this fixed support satisfies that condition, and enlarging the allowed
character order does not repair it.
\end{example}

\begin{theorem}[Exact minimum for the output condition; \stU]\label{spn:thm:minimum}
Among nontrivial finite $3$-power anticyclotomic characters of $K$ which are
trivial on the decomposition groups above $3$ and ramify only at primes lying
in split rational conjugate pairs, the minimum conductor norm is
$151^2=22801$. Exactly two characters attain this minimum. Both have order
three and conductor $(151)$.
\end{theorem}

\begin{proof}
For a split rational prime $q\ne3$, the minus local $3$-group is cyclic of
order $3^{h_q}$, where $h_q=v_3(q-1)$. If $h_q=0$, no such character can
ramify at $q$. Assume henceforth $h_q\ge1$.

Let $B_q$ be the minus $3$-ray group modulo $(q)$, and let $b$ be a minus
lift of Frobenius above $3$. Its image generates the ordinary class group.
Consequently $B_q$ is generated by its local group $R_q$ and $b$, with the
single additional class relation $3b=r_q\in R_q$. Up to inversion and
multiplication by a $3$-adic unit, $r_q$ is the $3$-primary residue of
$\beta/c(\beta)$, by \eqref{spn:eq:class-relations}. Killing Frobenius
therefore gives
\begin{equation}\label{spn:eq:ray-quotient}
 B_q/\langle b\rangle\simeq R_q/\langle r_q\rangle .
\end{equation}
It is nontrivial exactly when $\beta/c(\beta)$ is a cube modulo $q$. Choosing the conjugate root in the residue field only inverts the ratio
and does not change this test.

The exact finite calculation below includes every split prime $q\le151$ with
$q\equiv1\pmod3$. The root $t$ is a root of $T^2-T+6$ modulo $q$, and
\[
 u_q=\frac{1+2t}{3-2t}\pmod q,\qquad z_q=u_q^{(q-1)/3}\pmod q .
\]
\begin{center}
\begin{tabular}{rrrrrr}
\toprule
$q$ & $t$ & $u_q$ & $h_q$ & $z_q$ & $\abs{B_q/\langle b\rangle}$\\
\midrule
13  &5 &4  &1&9 &1\\
31  &8 &13 &1&5 &1\\
73  &30&54 &2&64&1\\
127 &28&78 &2&19&1\\
139 &17&102&1&42&1\\
151 &43&79 &1&1 &3\\
\bottomrule
\end{tabular}
\end{center}
At $151$ the quotient is $C_3$, giving exactly two nontrivial characters.
Their inertia is nontrivial, so the conductor is exactly $(151)$.

For minimality, unramified characters are killed by the condition at $3$
because the class of $\mathfrak s$ generates $\Cl(K)$. A character with one
rational prime in its support and smaller conductor would occur in an earlier
row and has just been excluded. A character with at least two distinct
rational support primes has their product at least $13\cdot31=403>151$.
Anticyclotomicity makes both conjugate primes ramify together. At residue
characteristic $q\ne3$, a finite $3$-power character is trivial on the higher
unit groups, which are pro-$q$; every ramified conductor exponent is one.
Thus a multi-prime support has conductor norm at least $403^2>151^2$. This
proves the claim without a squarefree-conductor assumption being added
separately.
\end{proof}

The table is reproduced by a certificate using only exact reduced-form
enumeration, norms, modular roots and modular
exponentiation.\footnote{\artifact{grh_transfer_bridges_2026-11-25/evidence/ray_character_checks.json}}
A two-prime alternative illustrates the role of the support. At $13$ and
$31$, both local groups are $C_3$ and both components of the relation $3b$
are nonzero. After changing their generators, killing $b$ imposes the
relation $(1,1)$ on $C_3^2$. The two nontrivial compatible characters have
local exponent vectors $(1,-1)$ and $(-1,1)$, and conductor norm
$(13\cdot31)^2=162409$; the one-prime support $151$ is cheaper.

\medskip
We now pass to the precursor condition and to small tame costs. Choose a
finite set $Q$ of primes of $K$ of prime norm, disjoint from $S$, with
$Q\cap cQ=\varnothing$. Thus each rational prime in the support splits in
$K$, and $Q$ contains one of its two primes. Consider nontrivial continuous
finite-order characters
\[
 \xi\colon G_K\longrightarrow\mu_{3^\infty},\qquad\psi=\xi/\xi^c,
\]
where $\xi$ is tamely ramified only on $Q$ and is trivial on the
decomposition groups at both members of $S$. Delete any primes of $Q$ where
$\xi$ is unramified, so that $Q$ is the actual support. Write
\begin{equation}\label{spn:eq:tame-cost-def}
 \spnTame(\psi)=\sum_{\mathfrak v\mid\mathfrak f(\psi)}
 v_3\bigl(1-\psi(\tau_{\mathfrak v})\bigr),\qquad v_3(3)=1,
\end{equation}
with $\tau_{\mathfrak v}$ a tame inertia generator; the quantity is
independent of those choices. If the local order at pair $i$ is $3^{n_i}$ and
the rational prime below it is $q_i$, then
\begin{equation}\label{spn:eq:source-twist-cost}
 \spnTame(\psi)=\sum_i3^{1-n_i},\qquad
 \Nm\mathfrak f(\psi)=\Bigl(\prod_iq_i\Bigr)^2 .
\end{equation}
Indeed, each primitive $3^{n_i}$th root contributes $1/(2\cdot3^{n_i-1})$ to
the valuation, and the pair contributes twice that number; the valuation
formula follows from $\Phi_{3^{n_i}}(1)=3$ and equality of the valuations of
its conjugate factors. Tameness makes each nontrivial conductor exponent
one.

Let
\[
 R_Q=\prod_{\mathfrak q\in Q}(\OO_K/\mathfrak q)^\times[3^\infty].
\]
The ordinary ray exact sequence has kernel $R_Q$ on the $3$-primary part,
because the units of $K$ are $\set{\pm1}$. Its quotient is $C_3$. A lift $b$
of $[\mathfrak s]$ satisfies $3b=r_\beta$, and the class of $c\mathfrak s$
is $r_3-b$, where $r_u$ denotes the local residue tuple of $u$, projected to
$R_Q$. A consistent inversion convention for principal ideals changes signs
but not the subgroups below. Consequently, imposing triviality on \emph{both}
decomposition groups above $3$ gives the exact quotient
\begin{equation}\label{spn:eq:source-twist-quotient}
 \frac{\Cl_Q(K)[3^\infty]}{\langle[\mathfrak s],[c\mathfrak s]\rangle}
 \simeq R_Q/\langle r_3,r_\beta\rangle .
\end{equation}
The characters of this quotient are exactly the allowed precursors $\xi$.
This is also the ray class sequence with the $S$-units quotiented out; the
class-field correspondence identifies it with the stated Galois-character
conditions \cite{Milne2020CFT}. For a single split rational prime $q$, put
$h=v_3(q-1)$ and choose a root $r$ of $r^2=-23$ in $\F_q$. The local group is
cyclic of order $3^h$. A precursor with local order at least $3^N$ exists if
and only if
\begin{equation}\label{spn:eq:residue-tests}
\begin{split}
 &q\equiv1\pmod{3^N},\qquad\legendre{-23}{q}=1,\\
 &3^{(q-1)/3^N}\equiv1\pmod q,\qquad(2+r)^{(q-1)/3^N}\equiv1\pmod q .
\end{split}
\end{equation}
These tests do not depend on which square root is chosen: the other value of
$\beta$ is $27/\beta$, and the test on $3$ already controls $27$.

\begin{theorem}[Exact conductor minima with the precursor condition; \stU]\label{spn:thm:source-twist-minimum}
For $1\le N\le5$, impose the budget $\spnTame(\psi)\le3^{1-N}$. Among all
characters $\psi=\xi/\xi^c$ just specified, including characters supported
on several split pairs, the exact minimum of $\Nm\mathfrak f(\psi)$ and the
number of output characters attaining it are as follows.
\[
\begin{array}{rcrrr}
\toprule
N & \text{Budget} & q_{\min} & \text{Minimum conductor norm} & \text{Number}\\
\midrule
1 & 1    & 151       & 22801             & 2\\
2 & 1/3  & 6481      & 42003361          & 24\\
3 & 1/9  & 6481      & 42003361          & 18\\
4 & 1/27 & 11767033  & 138463065623089   & 54\\
5 & 1/81 & 102906613 & 10589770999131769 & 162\\
\bottomrule
\end{array}
\]
Every minimum is attained on one conjugate pair, with conductor ideal
$(q_{\min})$ for $\psi$. For budget $1/3$ the minimum includes six
characters with local order $9$ and eighteen with local order $27$.
\end{theorem}

\begin{proof}
An exact search lists every prime in the progression $1\bmod3^N$ up to the
candidate bound, using a sieve of that progression and no probable-prime
tests, numerical approximations or class-group oracle, and applies
\eqref{spn:eq:residue-tests}; it finds the displayed first
primes.\footnote{\artifact{grh_ideal_densities_and_transfer_conditions_2026-11-26/data/source_twist_certificate.json}}
The exponents of the cyclic quotients in
\eqref{spn:eq:source-twist-quotient} at those primes are respectively
$1,3,3,4,5$.

This finite search also proves the minimum over arbitrary supports. If there
are at least two pairs and the budget is $3^{1-N}$, each $n_i$ must be at
least $N+1$: a single term with $n_i\le N$ would already exhaust or exceed
the budget, leaving no room for another positive term. Hence every support
prime is split and is $1\bmod3^{N+1}$. The first two such primes are
\[
\begin{array}{c|rrrrr}
N&1&2&3&4&5\\ \hline
q_1&73&163&163&487&42283\\
q_2&127&271&487&3889&52489.
\end{array}
\]
For $N=1,2,3,5$, their product exceeds the corresponding single-support
minimum $q_{\min}$, which excludes every multipair support.

For $N=4$, three pairs cost at least $487^3>11767033$ before the conductor is
squared. Since $3889^2>11767033$, every potentially cheaper two-pair support
must be $(487,q)$, with $q\in\set{3889,4861,8263,19441,19927}$. All these
primes have $v_3(q-1)=5$. Thus both local characters would have to have exact
order $3^5$. Choose generators of their local cyclic $3$-groups and reduce
the coordinate vectors of $(r_3,r_\beta)$ modulo $3$. One exact choice gives
\[
\begin{array}{c|rrrrrr}
q&487&3889&4861&8263&19441&19927\\ \hline
\log(3)\bmod3&1&0&2&1&0&0\\
\log(\beta)\bmod3&1&2&0&0&1&2.
\end{array}
\]
The column for $487$ is independent of every other column, even after
replacing the second row in that other column by its negative. The sign
change covers the other choice of the prime above $q$, since
$\log(\bar\beta)=3\log(3)-\log(\beta)$. Thus neither orientation admits
nonzero local character coefficients modulo $3$ satisfying both relations.
This excludes all five pairs.

Finally, at a cyclic quotient $C_{3^k}$, the number of characters with local
order at least $3^N$ is $\sum_{j=N}^{k}\varphi(3^j)$, which gives the final
column of the table. The map $\xi\mapsto\xi/\xi^c$ is injective for a fixed
half-support: its kernel would be unramified everywhere, while triviality at
$\mathfrak s$ kills every unramified character because $[\mathfrak s]$
generates the class group. Choosing the other half does not add new outputs,
since it replaces a precursor by the inverse of its conjugate.
\end{proof}

\begin{proposition}[The stronger condition changes the answer; \stU]\label{spn:prop:weak-minima}
If only the anticyclotomic output $\psi$ is required to be trivial above $3$,
the one-pair quotient is $C_{3^h}/\langle r_{\beta/\bar\beta}\rangle$, as in
\eqref{spn:eq:ray-quotient}, and the exact minimal support primes at budgets
$1,1/3,1/9,1/27,1/81,1/243$ are
\[
 151,\ 541,\ 6481,\ 94447,\ 238627,\ 1573183 .
\]
At budget $1/3$ the minimum conductor norm is $292681=541^2$ under the output
condition and $42003361=6481^2$ under the precursor condition.
\end{proposition}

\begin{proof}
The first entry is \cref{spn:thm:minimum}. The remaining minima, including
the exclusion of smaller multipair supports, are certified by the same exact
search.\footnote{\artifact{grh_ideal_densities_and_transfer_conditions_2026-11-26/data/source_twist_minima.csv}}
For the comparison at $q=541$, choose $r=40$, so $\beta=42$. Then
\[
 \Bigl(\frac{2+r}{2-r}\Bigr)^{60}\equiv1,\qquad3^{60}\equiv352,\qquad
 42^{60}\equiv129\pmod{541}.
\]
The weaker character exists at budget $1/3$, but the precursor condition
\eqref{spn:eq:residue-tests} fails.
\end{proof}

\begin{theorem}[A governing field for the precursor condition; \stU]\label{spn:thm:all-budgets}
For each $N\ge1$, set $k=3^N$ and
\[
 L_N=K\bigl(\zeta_k,3^{1/k},\beta^{1/k}\bigr).
\]
This is a finite Galois extension of $\Q$, of degree at most
$4\cdot3^{3N-1}$. Every rational prime $q$ splitting completely in $L_N$
supplies a precursor $\xi$ as above, with local order $3^N$ on one chosen
prime above $q$. Its output $\psi=\xi/\xi^c$ has conductor $(q)$ and
$\spnTame(\psi)=3^{1-N}$.

More generally, given an integer $m\ge1$ and $\eps>0$, choose
$N>\floor{\log_3m}$ such that $3^{1-N+\floor{\log_3m}}\le\eps$. Then a
character exists for which the first $m$ positive powers all have conductor
$(q)$ and satisfy $\spnTame(\psi^j)\le\eps$ for $1\le j\le m$. The support
can avoid any prescribed finite set of rational primes.
\end{theorem}

\begin{proof}
The field contains all conjugates of the displayed radicals. In particular,
$c\beta=27/\beta$, so a $k$th root of $c\beta$ is
$(3^{1/k})^3/\beta^{1/k}$, and the other roots differ by elements of
$\mu_k$. Thus $L_N/\Q$ is Galois. Adjoining the two radicals to $K(\zeta_k)$
gives
\[
 [L_N:\Q]\le2\varphi(k)k^2=4\cdot3^{3N-1}.
\]
Complete splitting makes $q$ split in $K$, gives $q\equiv1\pmod k$, and makes
$3$ and $\beta$ $k$th powers in the residue field. The tests
\eqref{spn:eq:residue-tests} therefore hold. The quotient in
\eqref{spn:eq:source-twist-quotient} has a character with exact local order
$k$, giving the claimed precursor. Chebotarev's theorem gives infinitely many
completely split primes, so any finite set may be excluded. This use is
unconditional and qualitative \cite{Milne2020CFT}.

For $1\le j\le m<3^N$, the inertia value of $\psi^j$ has order
$3^{N-v_3(j)}$. Hence
\[
 \spnTame(\psi^j)=3^{1-N+v_3(j)}\le3^{1-N+\floor{\log_3m}} .
\]
No such power loses ramification, proving the last assertions.
\end{proof}

For $0<\eps\le1$, choosing $N$ minimally in \cref{spn:thm:all-budgets} also
gives the explicit degree estimate $[L_N:\Q]<972\,(m/\eps)^3$. This bounds
the degree of the field governing the prime-selection problem. It does not
bound the least completely split prime or the conductor; in particular, a
least-prime estimate for rational arithmetic progressions cannot be applied
to $L_N$ merely from \cref{src:thm:dgrh}.

\begin{remark}[What these characters do and do not satisfy]\label{spn:rem:twist-scope}
The source already uses roots of $S$-units and Chebotarev to produce
auxiliary characters.\footnote{Declaration
\code{exists_anticyclotomicChar_of_frob_fixes_roots} of the source.} The
characters of \cref{spn:thm:source-twist-minimum,spn:thm:all-budgets} satisfy
the requirements isolated here: the local precursor condition, tame
ramification, triviality on the decomposition groups, and small valuation. They do not establish the simultaneous Selmer
length estimates, special-value bounds or prescribed Frobenius neighborhoods
needed for regular twists in the determinant argument (\cref{sec:sel}).
Also, $K=\Q(\sqrt{-23})$ is a benchmark base field, not the full cyclotomic
CM field of the source's admissible data; a base extension requires
recomputing its units and local conditions. No nonabelian Riemann hypothesis
follows from these constructions.
\end{remark}

%======================================================================

\subsection{The boundary for selection at a prescribed order}\label{cost:sec:boundary}
{\raggedright\sources{grh_method_transfers_2026-11-25_1; grh_exact_torsion_transfer_2026-11-25; grh-torsion-and-prime-sampling}}

All the positive statements above let the successful order vary in a window. The
flexibility cannot be replaced by a conclusion at one prescribed exact order with no
valuation loss, even for the comparison with the value at the trivial character.

Let $G=C_{p^m}$ with $m\ge2$, written additively with generator $g$ and identity $0$, and
let $\zeta$ be a primitive $p^m$-th root of unity. In the group ring
$\OO_{\Qp(\zeta)}[G]$ define
\[
 \vartheta=\prod_{a=1}^{p-1}\bigl([g]-\zeta^a[0]\bigr).
\]
A character $\chi$ evaluates $[g]$ at $\chi(g)$; the trivial character $\triv$ is the
augmentation.

\begin{proposition}[No maximal-order selection without loss; \stU]\label{cost:prop:maxorder}
One has
\[
 v(\vartheta(\triv))=(p-1)r_m=p^{1-m}=:t<1,\qquad v(\vartheta(\chi))\ge2t
\]
for every character $\chi$ of exact order $p^m$, where the second valuation may be
infinite. In particular no such character satisfies
$v(\vartheta(\chi))\le v(\vartheta(\triv))$.
\end{proposition}

\begin{proof}
Each $1-\zeta^a$, $1\le a\le p-1$, has valuation $r_m$, which gives the augmentation. Write
$\chi(g)=\zeta^b$ with $p\nmid b$. Exactly one $a\in\set{1,\ldots,p-1}$ satisfies
$a\equiv b\pmod p$. For every other $a$, $v(\zeta^b-\zeta^a)=r_m$. For the matching $a$ the
difference is zero or has valuation at least $p\,r_m$, because $p\mid b-a$. Summing gives
$v(\vartheta(\chi))\ge(p-2)r_m+p\,r_m=2t$.
\end{proof}

For $p=2$, $m=2$ this is $\vartheta=[g]-i[0]$: the augmentation has valuation $1/2$ and the
two primitive evaluations have valuations $+\infty$ and $1$. Here
$v_p(\abs{G[p]})=1$, so the weaker condition $v(\vartheta(\triv))<v_p(\abs{G[p]})$ would not
justify a maximal-order conclusion without loss. The threshold $p^{1-m}$ is exactly the one
appearing in the refined maximum-modulus principle of Kriz and
Nordentoft~\cite[Proposition~2.11]{KrizNordentoft2023Horizontal}, which in the cyclic case
gives the one-shell threshold relative to the trivial-character value; the example shows
that such ramification-sensitive thresholds cannot be dropped. The valuations are verified
exactly for $p=2,3,5,7$ and $m=2,3,4$ in
\artifact{grh_method_transfers_2026-11-25_1/data/torsion_checks.json}.

The example has a quantifier limit: $t$ tends to zero as $m$ grows, so it rules out no
loss, and any specified loss $\epsilon<t$, but by itself it does not rule out a statement
that first fixes $\epsilon>0$ and then takes $m$ large. The obstruction to a universal
exact-order selector for all polynomial inputs of bounded budget is the polynomial
$\Phi_{p^H}(1+T)$, which vanishes at every root of prescribed exact order $p^H$ while its
witness valuation at $0$ is $1$ (\cref{tors:rem:hypotheses}).

\subsection{Hypotheses that cannot be dropped}\label{tors:sec:limits}
\sources{grh_method_transfers_2026-11-25_1; grh_beyond_dirichlet_grh_2026-11-25; grh_exact_torsion_transfer_2026-11-25; grh_quantitative_twists_2026-11-25; grh_sharp_transfers_2026-11-25}

\begin{remark}\label{tors:rem:hypotheses}
(1) \emph{Bounded coefficients.} The series $\log(1+T)=\sum_{n\ge1}(-1)^{n+1}T^n/n$
converges for $v(T)>0$, has the nonzero value $\log(1+p^2)$ of valuation $2$, and
vanishes at every $\zeta-1$ with $\zeta$ of $p$-power order. Its coefficient valuations
$-\vp(n)$ are unbounded below, and no scalar multiple of it is integral. Likewise, for a
proposed $H$ and $x=p^R$ with $R\ge r$, the polynomial
$\bigl((1+T)^{p^H}-1\bigr)/\bigl((1+x)^{p^H}-1\bigr)$ has the value $1$ at $x$ and
vanishes at every $\zeta-1$ with $\zeta^{p^H}=1$; its denominator destroys integrality. A
family with a known common lower bound for the coefficient valuations can be scaled to an
integral family, but the scaling changes the budget.

(2) \emph{Positive radius.} For $F(T)=T$ one has $G_0(F)=0$, while $v(F(z))>0$ at every
torsion point. An exact comparison with the coefficient infimum $G_0$ is false.

(3) \emph{A total budget.} For $h\ge1$ let $F_h(T)=\Phi_{p^h}(1+T)\in\Zp[T]$ with the
witness $x_h=p^2$. Since $v\bigl((1+p^2)^{p^k}-1\bigr)=k+2$ for all $k\ge0$, including
$p=2$, one has $v(F_h(p^2))=1$ for every $h$. Each nontrivial root of unity of $p$-power
order annihilates one $F_h$. Infinitely many separate bounds $t_h\le1$ cannot replace a
finite total budget.

(4) \emph{The order cannot be prescribed.} The polynomial $\Phi_{p^H}(1+T)$ vanishes on
the whole shell $Z_H$, while its Gauss valuation at any radius $\rho\ge r_H$ is $1$. No
selection theorem under the present hypotheses can prescribe the exact order, and in
particular none can assert that the largest admissible order works. The theorems select
an order in an interval, and the dependence on the budget cannot be removed
(\cref{tors:rem:nonvanishing}).

(5) \emph{The moving target and separate exact bounds.} As recalled in
\cref{tors:sec:setup}, a bounded budget does not suffice for the exact comparison with
$G_{r_h}$ or for separate exact comparisons over $\OO$ with unbounded cardinality; see
\cref{sec:cost}.
\end{remark}

\paragraph{Relation to the literature.}
Kriz and Nordentoft develop horizontal $p$-adic $L$-functions, prove maximum-modulus
principles for Fourier transforms of measures on finite abelian groups, and deduce
simultaneous nonvanishing results for central values
\cite{KrizNordentoft2023Horizontal}; this is an arithmetic setting in which a selection
theorem of the present kind is a local ingredient. Serban studies the rigidity of the
intersection of torsion points with formal subschemes, including Lubin--Tate formal
groups \cite{Serban2018Infinitesimal,Serban2022Padic}; those results concern the structure
of torsion zero sets and do not give the finite budget estimates above by their stated
conclusions. Schefer uses nearest-root arguments and cyclotomic distance sums in
quantitative $p$-adic equidistribution of torsion points, with bounds involving the degree
\cite{Schefer2024Padic}. The least-occupied residue branch in the proof of
\cref{cost:lem:shell} belongs to the circle of ideas of Bhargava's $P$-orderings
\cite{Bhargava1997Porderings,Bhargava2009Porderings}. The statements proved here are the
explicit order intervals depending only on a Gauss budget at a positive radius, uniformly
in the degree. An application to an arithmetic family requires the interpolation
identity, the integral normalization and a nonzero witness with a known valuation, which
the selection theorems do not provide.

\subsection{Open questions}\label{tors:sec:open}
\sources{grh_quantitative_twists_2026-11-25; grh_fixed_field_transfer_2026-11-25; grh_beyond_dirichlet_grh_2026-11-25; grh_sharp_transfers_2026-11-25; grh_exact_torsion_transfer_2026-11-25; dirichlet_li_recovery}

\begin{question}[Families of fixed cardinality]\label{tors:q:cardinality}
For families of $n$ polynomials over $\OO$ with separate tolerances, determine the least
admissible exponent. For $n=1$ it is known exactly (\cref{sec:cost}), and for
unrestricted $n$ its logarithm is of order budget over tolerance
(\cref{tors:rem:tolerance}). In particular, can the exponent $n/(c\eta)$ in
\cref{tors:thm:fraction,tors:thm:packet} be lowered for $n\ge2$, as
\cref{tors:prop:linear} does for $n=1$?
\end{question}

\begin{question}[Dimension in the fixed-field theorem]\label{tors:q:dimension}
Determine the optimal dependence on the number $d$ of variables in
\cref{tors:thm:notower}. The factor $3^{d-1}$ and the range of slopes in
\cref{tors:lem:slopes} come from an interpolation estimate, not from an optimality
argument.
\end{question}

\begin{question}[Several powers]\label{tors:q:powers}
For $m=2$ and odd $p$ the endpoint $h_*(r/2)+\floor{2B}$ of \cref{tors:prop:powers}(a) is
optimal at positive half-integral budgets exactly when $\ord_p(2)$ is even
(\cref{tors:thm:twopowers,tors:rem:twopowers}). Is the endpoint $h_*(r/m)+\floor{mB}$
optimal for $m\ge3$, and what is the optimal endpoint when some of the powers $j\le m$ are
divisible by $p$, in particular for $p=2$ and $m=2$? More generally, when the covering
densities $c_h$ of \cref{tors:thm:cover} are not constant on $S$, is
$\sum_{h\in S}c_h$ the exact mixed-shell cost $C_s(S,J)$ for every $s\ge2$?
\end{question}

\begin{question}[Arithmetic families]\label{tors:q:arithmetic}
The extremal polynomials of \cref{sec:cost} are explicit, but an arbitrary polynomial
obstruction need not occur as a Selmer determinant or as an interpolation series.
Determine sharp exponents under restrictions that hold for actual arithmetic
interpolation families. For a specific horizontal interpolation system, can the local
bounds be combined with a witness and uniform estimates for the Euler factors to improve
a known arithmetic nonvanishing theorem?
\end{question}

\begin{question}[Density by conductor]\label{tors:q:density}
Can the shell density in \cref{tors:thm:fraction}, and the multiplicity in
\cref{tors:thm:arith}, be converted into a density statement when conductors are ordered
by norm? This requires a bound for the norms of the auxiliary primes in $Q$ and control
of the variation of the parameter space of the characters.
\end{question}

\begin{question}[Which objective does each use of the selector need?]\label{tors:q:objective}
The denominator bookkeeping of the source often adds losses, which is the aggregate
comparison of \cref{tors:cor:packets}(a) with its optimal exponent; other steps impose a
separate inequality for each Selmer datum, for which the aggregate theorem is
insufficient. Determine, for each use of the selector in the source, which comparison is
needed, and hence where the exponent of \cref{tors:thm:selection} or of
\cref{tors:thm:fixedfield} can replace the harmonic one.
\end{question}

\subsection{Open questions}\label{cost:sec:open}
{\raggedright\sources{grh-torsion-and-prime-sampling; grh-torsion-and-zero-data; grh_method_transfers_sharp_torsion; grh_sharp_torsion_and_symplectic_blocks; grh_fixed_field_transfer_2026-11-25; grh_method_transfers_research; grh_fixed_cardinality_torsion}}

\begin{question}[Fixed packet cardinality]\label{cost:q:cardinality}
At tolerance zero the cost of blocking with $n$ members is $s_{h_0}\log_p^{\circ(n-1)}H$
up to the error terms of Theorem~\ref{cost:thm:cardinality}, and every finite instance is
computed by Theorem~\ref{cost:thm:cardfinite}.
\begin{enumerate}[label=\textup{(\alph*)}]
\item Can the upper error $O(\log B)$ in \eqref{cost:eq:tower} be reduced to $O(1)$, or the
next term be determined for $n=2$? Is the optimum always attained by consecutive blocks with
one member each, or do mixed assignments within a shell do better?
\item At a fixed positive tolerance $\delta$, what is the asymptotic of $H_n(B,\delta)$ as
$B\to\infty$, and what is the joint transition when $\delta$ decreases with $B$? The tower
law at tolerance zero does not determine it.
\item Over a fixed finite coefficient field, where the root multisets must be Galois
stable, how do ramification and cyclotomic intersections constrain the activations and
change the height of the tower?
\item Does the delayed product potential extend to the settings \textup{(LT)} and
\textup{(H)}, and to formal groups over ramified bases with overlapping torsion valuations
(\cref{cost:sec:ramified})? The cyclotomic proof uses the equal spacing ratios
$r_{h-1}=pr_h$ and the exact geometry of the shells.
\end{enumerate}
\end{question}

\begin{question}[Positive tolerance over a general coefficient field]\label{cost:q:field}
For a finite extension $K/\Qp$ that is neither abelian nor of the form $\Qp(\pi)$ with
$\pi^p=pu$ and $p$ odd, find a closed formula for the shell costs $c_{K,h}(\delta)$ at
positive tolerance. Which ramification invariants of $K$ determine them, as the quadratic
defect of $-1$ does at order four (Proposition~\ref{cost:prop:fourth})? By Remark~\ref{cost:rem:twofields}
the value group and the cyclotomic intersection do not suffice.
Proposition~\ref{cost:prop:mod} supplies a finite test against which candidate formulas can be checked.
\end{question}

\begin{question}[One polynomial over a prescribed field]\label{cost:q:single}
For a fixed ramified coefficient field $K$, which polynomials over $K$ attain the
unrestricted threshold of Theorem~\ref{cost:thm:minimax}, and what is the closed form of the
recursion of Proposition~\ref{cost:prop:single}? Likewise, if the polynomial in Theorem~\ref{cost:thm:finite}
must have coefficients in the field of definition of the formal group, as for the elliptic
curve of Example~\ref{cost:ex:elliptic}, which of the exact thresholds can still be attained? The
unrestricted construction selects individual torsion points even when their Galois
conjugates are omitted.
\end{question}

\begin{question}[Several variables at positive tolerance]\label{cost:q:several}
For several coordinates and total positive excess at tolerance $\delta>0$, is the leading
exponential constant of Corollary~\ref{cost:cor:depth} attained by an obstruction in one variable?
Sequential substitution gives upper bounds, but it can spend the same tolerance at
successive coordinates.
\end{question}
